\documentclass[11pt,a4paper]{article}

\usepackage[utf8]{inputenc}
\usepackage[T1]{fontenc}
\usepackage{lmodern}
\usepackage[margin=1in]{geometry}
\usepackage{amsmath,amssymb,amsthm,mathtools}
\usepackage{booktabs}
\usepackage{array}
\usepackage{enumitem}
\usepackage{xcolor}
\usepackage{hyperref}

\DeclareUnicodeCharacter{03B7}{\ensuremath{\eta}}
\DeclareUnicodeCharacter{03BC}{\ensuremath{\mu}}
\DeclareUnicodeCharacter{03C4}{\ensuremath{\tau}}
\DeclareUnicodeCharacter{2192}{\ensuremath{\to}}

\hypersetup{
  colorlinks=true,
  linkcolor=blue!60!black,
  urlcolor=blue!60!black,
  citecolor=blue!60!black,
  pdftitle={The CoEndoYoneda Lemma and its Relationship to a Pointfree Categorical Characterization of Purity},
  pdfauthor={Luc Duponcheel}
}

\newtheorem{theorem}{Theorem}[section]
\newtheorem{lemma}[theorem]{Lemma}
\newtheorem{proposition}[theorem]{Proposition}
\newtheorem{corollary}[theorem]{Corollary}
\theoremstyle{definition}
\newtheorem{definition}[theorem]{Definition}
\newtheorem{example}[theorem]{Example}
\theoremstyle{remark}
\newtheorem{remark}[theorem]{Remark}

\newcommand{\Set}{\mathbf{SET}}
\newcommand{\Hom}{\mathrm{Hom}}
\newcommand{\NatTrans}{\mathrm{NatTrans}}
\newcommand{\Id}{\mathrm{Id}}
\newcommand{\id}{\mathrm{id}}
\newcommand{\CYF}{\mathrm{CYF}}
\newcommand{\CYEF}{\mathrm{CYEF}}
\newcommand{\FF}{\mathrm{FF}}
\newcommand{\KFF}{\mathrm{KFF}}
\newcommand{\EF}{\mathrm{EF}}
\newcommand{\K}{\mathbf{K}}
\newcommand{\ttstar}{\mathrm{tt}^*}
\newcommand{\const}{\mathrm{const}}
\newcommand{\starK}{\mathbin{\star_{\K}}}
\newcommand{\iotaZ}{\iota_Z}

\setlength{\parindent}{0pt}
\setlength{\parskip}{6pt plus 2pt minus 1pt}
\setlength{\emergencystretch}{3em}

\title{\textbf{
  A CoEndoYoneda Lemma and its Relationship to\\
  a Pointfree Categorical Characterization of Purity
}}

\author{Luc Duponcheel}

\date{November 2026}

\begin{document}

\maketitle

\begin{abstract}
\noindent The classical CoYoneda lemma is a pointful statement about $\Set$-valued \emph{functors}
$\CYF(Z) = \Hom_{\mathcal{C}}(Z, -) : \mathcal{C} \to \Set$. In the special case where the base
category is $\mathcal{C} = \Set$ itself, $\CYF(Z) = \Hom_{\Set}(Z, -) : \Set \to \Set$ is an
\emph{endofunctor}.

\vspace{4pt}

\noindent Motivated by pointfree functional programming, where programs are morphisms in a category
$\mathcal{C}$ equipped with a \emph{functional functor} $\FF : \Set \to \mathcal{C}$ lifting pure
functions to morphisms, pure programs, one can compose $\CYF(Z)$ with $\FF$  to obtain an
\emph{endofunctor} $\CYEF(Z) = \FF \circ \CYF(Z) : \mathcal{C} \to \mathcal{C}$. 

\vspace{4pt}

\noindent We introduce \emph{Basic Functional Categories} and \emph{Functional Categories} relative
to $\FF$, equipped with a global-like endofunctor $G = \CYEF(\FF(T))$, where $T$ is the terminal
set, a global-like element generator function $g_Z : Z \to \Hom_{\mathcal{C}}(\FF(T), \FF(Z))$ in
$\Set$ defined by $g_Z(z) = \FF(\const_z)$ for the "constant $z$" function $\const_z : T \to Z$, a
unit-like natural transformation $\eta : \Id_{\mathcal{C}} \Rightarrow G$ satisfying the parametric
equational law $\eta_{\FF(Z)} = \FF(g_Z)$ (Equation~\eqref{eq:eta-eq}), and a multiplication-like
natural transformation $\mu : G \circ G \Rightarrow G$ satisfying a left unit law
$\eta_{G(Z)} \star \mu_Z = \id_{G(Z)}$ (Equation~\eqref{eq:id-left}). 

\vspace{4pt}

\noindent We then prove the \emph{CoEndoYoneda Isomorphism Theorem}
(Theorem~\ref{thm:coendoyoneda}): for every object $Z \in \mathrm{ob}(\mathcal{C})$ and endofunctor
$\EF : \mathcal{C} \to \mathcal{C}$, there is a canonical pointfree bijection
\[
  \NatTrans\bigl(\CYEF(Z),\, G \circ \EF\bigr)
  \;\cong\;
  \Hom_{\mathcal{C}}\bigl(\FF(T),\, (G \circ \EF)(Z)\bigr),
\]
which generalizes the classical CoYoneda canonical pointful bijection
\[
  \NatTrans\bigl(\CYF(Z),\, EF\bigr) 
  \;\cong\; 
  EF(Z)
\]
for endofunctors on $\Set$ (recovered when
$\mathcal{C} = \Set$ and $\FF = \Id_{\Set}$, Proposition~\ref{prop:set-fc} and
Subsection~\ref{subsec:coyoneda-comparison}) to any Functional Category $(\mathcal{C}, \FF)$.

\vspace{4pt}

\noindent Evaluating the forward isomorphism at $\EF = \Id_{\mathcal{C}}$ and $Z = \FF(T)$ on the
identity transformation $\id_G$ recovers the terminal unit component
$\eta_{\FF(T)} = g_{\FF(T)}(\id_{\FF(T)}) = g_{\FF(T)}(g_T(\ttstar))$, where $\ttstar$ is the single
value of Terminal type $T$ (Proposition~\ref{prop:origin-eta}).

\vspace{4pt}

\noindent Next, we investigate which effectful categories of programs satisfy the Basic Functional
Category axioms. For the Kleisli category $\K(M)$ of a monad $(M, \mathrm{mPure}, \mathrm{mApply})$
on $\Set$ equipped with the canonical identity-on-objects embedding $\KFF : \Set \to \K(M)$, we prove
our second main result (Theorem~\ref{thm:kleisli-purity}): \emph{$(\K(M), \KFF)$ is a Basic Functional
Category if and only if the monad $M$ is idempotent} ($\mathrm{mMap}\;\mathrm{mPure} = \mathrm{mPure}$,
equivalently $\mathrm{mApply}\;(\mathrm{mPure} \circ \mathrm{mPure}) = \mathrm{mPure}$). Remarkably,
even the specialized terminal unit law $\eta_{\FF(T)} = g_{\FF(T)}(g_T(\ttstar))$
(Proposition~\ref{prop:eta-eq-T}) suffices to force global idempotency of $M$, which in turn yields the
full parametric law~\eqref{eq:eta-eq} definitionally. 

\vspace{4pt}

\noindent Thus, the Basic Functional Category law provides an elementary, pointfree,
specification-level characterization of \emph{computational purity}. We further derive the
terminal subsingleton theorem ($M(T) \cong T$ for every pure monad,
Theorem~\ref{thm:terminal-subsingleton}), provide complete equational classifications of seven
standard computational effects ($\mathrm{Option}$, $\mathrm{List}$, $\mathrm{Read}\;R$,
$\mathrm{Write}\;W$, $\mathrm{State}\;S$, $\mathrm{Cont}\;R$, and $\mathrm{ErrorState}\;E\;S$,
Theorem~\ref{thm:examples}), prove a strict separation theorem (Theorem~\ref{thm:strict-separation})
with an explicit counterexample showing that $\mathrm{Read}\;R$ ($|R| \ge 2$) satisfies
$M(T) \cong T$ while failing the Basic Functional Category condition, and discuss the scope,
implications for commutative effects, and limitations of the framework
(Subsection~\ref{subsec:limitations}).

\vspace{4pt}

\noindent \textbf{Provenance and Attribution Statement.}
The mathematical formulation of Basic Functional Categories, Functional Categories, the CoEndoYoneda
Lemma, and the equational proof proper of the CoEndoYoneda Isomorphism Theorem were conceived and
proved by the author (Luc Duponcheel)~\cite{DuponcheelAcademia}. By contrast, the Lean~4
mechanization and the discovery and proof of the exact equivalence between the Basic Functional
Category structure on $\K(M)$ and monad idempotency/purity were established collaboratively
using \emph{Gemini} (Google DeepMind Antigravity) guided by mathematical prompts from the author. Both
the Cubical Agda development~\cite{CubicalAgdaRepo} and the Palomar-verified Lean~4
formalization~\cite{PalomarLeanRepo} are machine-checked with zero unproved axioms or
postulates.
\end{abstract}

\noindent \textbf{Keywords:} CoYoneda lemma, pointfree category theory, functional categories, Kleisli
categories, idempotent monads, computational purity, Cubical Agda, Lean~4.

\noindent \textbf{MSC 2020:} 18A40, 18C15, 18C20, 68N18, 68V20. \quad
\textbf{arXiv:} \texttt{cs.PL}, \texttt{math.CT}, \texttt{cs.LO}.

\newpage
\tableofcontents
\newpage

\section{Introduction and Motivation}\label{sec:intro}

\subsection{From Pointful CoYoneda to Pointfree CoEndoYoneda}

In classical category theory, CoYoneda theorems and proofs deal with $\Set$-valued functors (co-presheaves)
\[
  \CYF(Z) : \mathcal{C} \longrightarrow \Set,
  \qquad
  Y \longmapsto \Hom_{\mathcal{C}}(Z, Y),
  \qquad
  (y2x : Y \to X) \longmapsto (z2y \mapsto z2y \star y2x),
\]
where throughout this paper we write $f \star h : Z \to X$ for the diagrammatic composition of
$f : Z \to Y$ followed by $h : Y \to X$ in $\mathcal{C}$. For any functor $F : \mathcal{C} \to \Set$,
the classical CoYoneda lemma establishes a natural bijection
\begin{equation}\label{eq:classical-coyoneda}
  \NatTrans\bigl(\CYF(Z),\, F\bigr) \;\cong\; F(Z),
\end{equation}
mediated by the universal element $\id_Z \in \Hom_{\mathcal{C}}(Z, Z) = \CYF(Z)(Z)$ via the forward map
$\tau \mapsto \tau_Z(\id_Z)$ and the backward map $fz \mapsto \bigl(Y \mapsto (z2y \mapsto F(z2y)(fz))\bigr)$.

Notice an important special case of~\eqref{eq:classical-coyoneda}: \emph{when the base category
$\mathcal{C}$ is $\Set$ itself, the CoYoneda functor $\CYF(Z) = \Hom_{\Set}(Z, -) : \Set \to \Set$ and
any target functor $F = \EF : \Set \to \Set$ are both \textbf{endofunctors} on $\Set$}. Furthermore, in
$\Set$, the set of elements $\EF(Z)$ is canonically isomorphic to the hom-set of global elements
$\Hom_{\Set}(T, \EF(Z)) = (T \to \EF(Z))$ out of the terminal singleton set $T = \{\ttstar\}$. Thus, on
$\mathcal{C} = \Set$, the classical CoYoneda lemma~\eqref{eq:classical-coyoneda} is already an
\emph{endofunctor} statement relating natural transformations between endofunctors on $\Set$ to
morphisms out of $T$:
\[
  \NatTrans\bigl(\CYF(Z),\, \EF\bigr)
  \;\cong\; \EF(Z)
  \;\cong\; \Hom_{\Set}\bigl(T,\, \EF(Z)\bigr).
\]
For a general category $\mathcal{C}$, however, classical CoYoneda theorems are inherently
\emph{pointful}: they study objects and functors by passing outside $\mathcal{C}$ into the category
$\Set$ of hom-sets and set-theoretic elements. At the same time, they serve as the bridge from a
pointful approach to a \emph{pointfree} approach, characterizing objects by the morphisms out of them.

Motivated by \emph{pointfree functional programming}, such as my Program Specification Based
Programming (\texttt{PSBP}) architecture~\cite{PSBP}, where pure functions can be lifted to program
morphisms, a natural foundational question arises:
\begin{quote}
\emph{What if the category $\mathcal{C}$ is also equipped with a ``functional functor''
$\FF : \Set \to \mathcal{C}$ mapping functions to morphisms?}
\end{quote}
Whenever such a functor $\FF : \Set \to \mathcal{C}$ exists, every $\Set$-valued CoYoneda functor
$\CYF(Z) : \mathcal{C} \to \Set$ can be composed with $\FF$ to form an \emph{endofunctor} on
$\mathcal{C}$, just as in the $\mathcal{C} = \Set$ case:
\[
  \CYEF(Z) \;\coloneqq\; \FF \circ \CYF(Z) \;:\; \mathcal{C} \longrightarrow \mathcal{C}.
\]
We can then ask: 
\begin{quote}
\emph{Are there pointfree \textbf{CoEndoYoneda} theorems that generalize the
$\mathcal{C} = \Set$ endofunctor CoYoneda lemma so that they can be formulated and proved, in a
pointfree way, entirely in terms of endofunctors on $\mathcal{C}$?}
\end{quote}

In Sections~\ref{sec:bfc}--\ref{sec:coendoyoneda}, we show that this is indeed possible for categories
equipped with $\FF : \Set \to \mathcal{C}$ together with two natural transformations ($\eta$ and $\mu$)
and two simple equational laws (Equations~\eqref{eq:eta-eq} and~\eqref{eq:id-left}). 

In Subsection~\ref{subsec:coyoneda-comparison}, we provide a detailed comparison of how the
CoEndoYoneda isomorphism relates to the classical CoYoneda lemma in terms of naturality,
universality, generalization, and specialization.

\subsection{Theory Meets Practice: ``Every Disadvantage Has Its Advantage''}

Having established the CoEndoYoneda theorem, the natural practical question is: 
\begin{quote}
\emph{For which
categories of effectful programs can we apply this theorem?}
\end{quote}

Originally, I hoped that the CoEndoYoneda framework could be instantiated for at least a
broad class of effectful (monadic) computations. However, testing the axioms on Kleisli categories
$\K(M)$ revealed a surprise: the parametric unit law~\eqref{eq:eta-eq} (and even its terminal
specialization~\eqref{eq:eta-eq-T} in Proposition~\ref{prop:eta-eq-T}) fails for impure effects such
as $\mathrm{Option}$, $\mathrm{List}$, $\mathrm{State}\;S$, and $\mathrm{ErrorState}\;E\;S$. At first
sight, this appeared to be bad news: the CoEndoYoneda bijection cannot be used for arbitrary effectful
computations.

Yet, in the spirit of Willem van Hanegem's famous Dutch aphorism \emph{``elk nadeel heb se voordeel''}
(\emph{``every disadvantage has its advantage''}, often misattributed to Johan Cruyff), this apparent
limitation turned out to be a theoretical gift:
\begin{quote}
\emph{For a Kleisli category $\K(M)$ with its canonical identity-on-objects functional functor
$\KFF : \Set \to \K(M)$, the Basic Functional Category structure is \textbf{equivalent} to the
idempotency (purity) of the monad $M$ (Theorem~\ref{thm:kleisli-purity}).}
\end{quote}

While monads are inherently \emph{pointful}, our Basic Functional Category axiom provides an
elementary, \emph{pointfree specification-level characterization of purity} directly on the category
of programs.

\subsection{Related Work on Categorical and Pointfree Characterizations of Purity}\label{sec:related}

To position our contribution within the programming languages and category theory literature, we
compare our Basic Functional Category characterization against established categorical frameworks for
distinguishing pure and effectful computations:

\begin{enumerate}[label=(\arabic*),nosep]
  \item \textbf{Moggi's computational $\lambda$-calculus and strong monads~\cite{Moggi1991}:}
  Moggi pioneered the modeling of computational effects via strong monads $(M, \eta, \mu, \mathrm{st})$
  on a cartesian closed category $\mathcal{C}$. In Moggi's setting, pure programs are identified
  extrinsically as those Kleisli morphisms in $\K(M)$ that factor through the canonical inclusion
  functor $\KFF : \mathcal{C} \to \K(M)$, whereas our Theorem~\ref{thm:kleisli-purity} characterizes
  when the entire Kleisli category $\K(M)$ collapses to pure computation via the existence of the
  CoYoneda-induced natural transformation $\eta : \Id_{\K(M)} \Rightarrow G$.
\vspace{4pt}
  \item \textbf{Premonoidal categories, Freyd categories, and central morphisms~\cite{PowerRobinson1997,PowerThielecke1999,LevyPowerThielecke2003}:}
  Power, Robinson, and Thielecke (and Levy in Fine-Grain Call-by-Value) axiomatize effectful
  categories directly without mentioning a monad via \emph{premonoidal categories} and \emph{Freyd
  categories} (an identity-on-objects symmetric premonoidal functor $J : \mathcal{V} \to \mathcal{K}$
  from a cartesian category of values $\mathcal{V}$ to a premonoidal category of computations
  $\mathcal{K}$). In a premonoidal category, a morphism $f : X \to Y$ is characterized pointfreely as
  \emph{central} (effect-commuting) if $(f \ltimes g) = (f \rtimes g)$ for all morphisms $g$. While
  centrality captures commutativity of effects (and hence purity relative to non-commuting effects),
  it requires symmetric premonoidal tensor structure on $\mathcal{K}$ and still admits commutative
  effects (such as $\mathrm{Read}\;R$ or commutative $\mathrm{Write}\;W$) in the center. By contrast,
  our Basic Functional Category condition requires no monoidal or premonoidal tensor structure on
  $\mathcal{C}$, and eliminates even commutative effects such as $\mathrm{Read}\;R$ (for $|R| \ge 2$)
  and commutative $\mathrm{Write}\;W$ (for $|W| \ge 2$).
\vspace{4pt}
  \item \textbf{Führmann's direct models of the computational $\lambda$-calculus and thunkability~\cite{Fuhrmann1999}:}
  Führmann gave a pointfree categorical semantics of call-by-value programming by equipping a
  category $\mathcal{K}$ with a \emph{thunking endofunctor} $T : \mathcal{K} \to \mathcal{K}$ carrying
  a comonad-like structure, characterizing pure (``thunkable'') morphisms $f : X \to Y$ pointfreely by
  commutation with thunking ($f \star \vartheta_Y = \vartheta_X \star T(f)$, copyability, and
  discardability). Our global-like endofunctor $G = \FF \circ \Hom_{\mathcal{C}}(\FF(T), -)$ plays a
  role analogous to an intrinsic thunking functor $Y \mapsto (T \to Y)$ constructed canonically from
  $\FF$ and $\CYF(\FF(T))$. However, whereas Führmann axiomatizes thunkable morphisms inside a general
  effectful category using comonadic structural laws, our Definition~\ref{def:bfc} requires only a
  single natural transformation $\eta : \Id_{\mathcal{C}} \Rightarrow G$ and the single equation
  $\eta_{\FF(Z)} = \FF(g)$ (or $\eta_{\FF(T)} = g(g(\ttstar))$).
\vspace{4pt}
  \item \textbf{Lawvere theories and algebraic effects~\cite{Lawvere1963,PlotkinPower2002,PlotkinPower2003}:}
  In the algebraic theory of effects developed by Plotkin and Power, finitary monads on $\Set$
  correspond to Lawvere theories $\mathcal{L}$, and effectful operations are modeled by generic effects
  or algebraic operations. Within this framework, idempotent monads on $\Set$ correspond to
  \emph{subterminal} (or propositional) Lawvere theories in which every $n$-ary operation reduces to a
  projection or constant, and on general categories idempotent monads correspond to reflective
  subcategories and well-pointed endofunctors ($M\eta = \eta M$)~\cite{Kelly1980,Borceux1994,nLabIdempotentMonad,nLabWellPointed}.
  Our Theorem~\ref{thm:kleisli-purity} bridges this classical monadic characterization ($M\eta = \eta M$)
  with the pointfree CoEndoYoneda structure on $\K(M)$.
\end{enumerate}

\subsection{Provenance and Attribution of Results}\label{sec:provenance}

In accordance with transparent human--AI scientific attribution standards (aiXiv and the Palomar
Registry), we state the exact provenance of each contribution in this work:
\begin{enumerate}[label=(\arabic*),nosep]
  \item \textbf{Human-authored core theory and proof proper:} The categorical definitions of Basic
  Functional Categories and Functional Categories, the statement of the CoEndoYoneda Isomorphism
  Theorem (Theorem~\ref{thm:coendoyoneda}), and its complete pointfree equational proof proper were
  conceived, designed, and proved by Luc Duponcheel~\cite{DuponcheelAcademia}.
\vspace{4pt}
  \item \textbf{AI-assisted purity equivalence and formal mechanizations:} The Lean~4 formalization
  and the discovery and proof of the exact equivalence between the Basic Functional Category structure
  on $\K(M)$ and monad idempotency/purity (Theorem~\ref{thm:kleisli-purity}), the
  terminal subsingleton diagnostic (Theorem~\ref{thm:terminal-subsingleton}), the formal classification
  of the seven computational effects and strict separation theorem (Section~\ref{sec:examples}), and the
  synchronized Cubical Agda and Lean~4 implementations were developed using \textbf{Gemini} (Google
  DeepMind Antigravity) based upon targeted mathematical prompts and iterative guidance from the author.
\end{enumerate}

\section{Basic Functional Categories}\label{sec:bfc}

\paragraph{Conventions and Notation.}
Throughout this paper, let $\mathcal{C}$ be a locally small category at universe level $\ell$ with object
collection $\mathrm{ob}(\mathcal{C})$, and for objects $Z, Y \in \mathrm{ob}(\mathcal{C})$, let
$\Hom_{\mathcal{C}}(Z, Y) \in \mathrm{ob}(\Set_\ell)$ denote the hom-set of morphisms in $\mathcal{C}$
from $Z$ to $Y$. Because $\mathcal{C}$ is locally small at level $\ell$, every hom-set
$\Hom_{\mathcal{C}}(Z, Y)$ is an object of the category of sets $\Set_\ell$, and a morphism $f$ from $Z$
to $Y$ in $\mathcal{C}$ is an element of that hom-set: we write $f : Z \to Y$ in $\mathcal{C}$ and
$f \in \Hom_{\mathcal{C}}(Z, Y)$ synonymously.\footnote{In dependent type-theoretic proof assistants such
as Cubical Agda and Lean~4, the hom-set $\Hom_{\mathcal{C}}(Z, Y)$ is represented as a type in
$\mathrm{Type}\;\ell$, so set-theoretic membership $f \in \Hom_{\mathcal{C}}(Z, Y)$ is expressed by the
typing judgment $f : \Hom_{\mathcal{C}}(Z, Y)$. In this paper we follow standard categorical notation
$f : Z \to Y$ in $\mathcal{C}$ (or $f \in \Hom_{\mathcal{C}}(Z, Y)$) to keep the distinction between an
individual morphism $f$ and the hom-set $\Hom_{\mathcal{C}}(Z, Y) \in \mathrm{ob}(\Set_\ell)$ completely
explicit.}
We write $\id_Z : Z \to Z$ for the identity morphism on $Z$, and $f \star h : Z \to X$ for the
\emph{diagrammatic composition} of $f : Z \to Y$ followed by $h : Y \to X$ in an abstract category
$\mathcal{C}$ (for Kleisli categories $\K(M)$ in Section~\ref{sec:purity}, we write $\starK$ for Kleisli
composition to distinguish it from ordinary function composition $\circ$ in $\Set_\ell$). For a functor
$F : \mathcal{C} \to \mathcal{D}$, we write $F(Z) \in \mathrm{ob}(\mathcal{D})$ for its action on an
object $Z \in \mathrm{ob}(\mathcal{C})$ and $F(f) : F(Z) \to F(Y)$ for its action on a morphism
$f : Z \to Y$. Given functors $F, H : \mathcal{C} \to \mathcal{D}$, a natural transformation
$\tau : F \Rightarrow H$ (an element of $\NatTrans(F, H)$) assigns to each object
$Y \in \mathrm{ob}(\mathcal{C})$ a component morphism $\tau_Y : F(Y) \to H(Y)$ in $\mathcal{D}$ such
that $F(f) \star \tau_X = \tau_Y \star H(f)$ for every morphism $f : Y \to X$ in $\mathcal{C}$.
Finally, $\Set_\ell$ (abbreviated $\Set$ when the universe level $\ell$ is clear from context) denotes
the category of sets at universe level $\ell$ (hSets in Cubical Agda, $\mathrm{Type}\;u$ in Lean~4), in
which objects are sets and morphisms $f : Z \to Y$ are functions.

\begin{definition}\label{def:cyef}
Given a locally small category $\mathcal{C}$ and a \emph{functional functor}
$\FF : \Set_\ell \to \mathcal{C}$, we define:
\begin{enumerate}[label=(\roman*),nosep]
  \item The \emph{CoYoneda embedding} $\CYF : \mathrm{ob}(\mathcal{C}) \to \mathrm{Functor}(\mathcal{C}, \Set_\ell)$:
  for each $Z \in \mathrm{ob}(\mathcal{C})$,
  \begin{align*}
    \CYF(Z)(Y) &\coloneqq \Hom_{\mathcal{C}}(Z, Y) \;\in\; \mathrm{ob}(\Set_\ell), \\
    \CYF(Z)(y2x) &\coloneqq \bigl(z2y \mapsto z2y \star y2x\bigr)
      \;:\; \Hom_{\mathcal{C}}(Z, Y) \to \Hom_{\mathcal{C}}(Z, X).
  \end{align*}
\vspace{4 pt}
  \item The \emph{CoYoneda endofunctor} $\CYEF : \mathrm{ob}(\mathcal{C}) \to \mathrm{Functor}(\mathcal{C}, \mathcal{C})$:
  \[
    \CYEF(Z) \;\coloneqq\; \FF \circ \CYF(Z) \;:\; \mathcal{C} \longrightarrow \mathcal{C}.
  \]
\vspace{4 pt}
  \item The \emph{terminal set} $T \coloneqq \{ \ttstar \} \in \mathrm{ob}(\Set_\ell)$, a singleton set
  with unique element $\ttstar \in T$.
\vspace{4 pt}
  \item The \emph{global-like endofunctor} $G : \mathcal{C} \to \mathcal{C}$:
  \[
    G \;\coloneqq\; \CYEF\bigl(\FF(T)\bigr) \;=\; \FF \circ \CYF\bigl(\FF(T)\bigr),
  \]
  so that $G(Y) = \FF\bigl(\Hom_{\mathcal{C}}(\FF(T), Y)\bigr)$ for $Y \in \mathrm{ob}(\mathcal{C})$ and
  $G(y2x) = \FF(f \mapsto f \star y2x) : G(Y) \to G(X)$ for $y2x : Y \to X$.
\vspace{4 pt}
  \item The \emph{global-element generator} (a family of functions in $\Set_\ell$ indexed by sets
  $Z \in \mathrm{ob}(\Set_\ell)$): for each set $Z \in \mathrm{ob}(\Set_\ell)$ and each element
  $z \in Z$, let $\const_z : T \to Z$ denote the constant function $\const_z(t) \coloneqq z$ in
  $\Set_\ell$. Applying the functor $\FF : \Set_\ell \to \mathcal{C}$ to the morphism
  $\const_z : T \to Z$ in $\Set_\ell$ yields a morphism $\FF(\const_z) : \FF(T) \to \FF(Z)$ in
  $\mathcal{C}$, i.e., an element of the hom-set $\Hom_{\mathcal{C}}(\FF(T), \FF(Z)) \in \mathrm{ob}(\Set_\ell)$.
  We therefore define the function $g_Z : Z \to \Hom_{\mathcal{C}}(\FF(T), \FF(Z))$ in $\Set_\ell$
  (often written simply $g$ when the domain set $Z$ is clear from context) by:
  \begin{equation}\label{eq:def-g}
    \begin{aligned}
      g_Z &\;:\; Z \longrightarrow \Hom_{\mathcal{C}}\bigl(\FF(T),\, \FF(Z)\bigr)
      \quad\text{in }\Set_\ell, \\
      g_Z(z) &\;\coloneqq\; \FF(\const_z) \;\in\; \Hom_{\mathcal{C}}\bigl(\FF(T),\, \FF(Z)\bigr).
    \end{aligned}
  \end{equation}
\end{enumerate}
\end{definition}

\begin{remark}[Two-level typing of $g_Z$ and $\FF(g_Z)$]\label{rem:g-typing}
Because the construction of $g_Z$ and $\FF(g_Z)$ crosses back and forth between $\Set_\ell$ and
$\mathcal{C}$ via the functor $\FF : \Set_\ell \to \mathcal{C}$, it is helpful to trace the exact types
at each step:
\begin{enumerate}[label=(\arabic*),nosep]
  \item For each element $z \in Z$, $\const_z = (\_ \mapsto z) : T \to Z$ is a function (a morphism in
  $\Set_\ell$).
\vspace{4 pt}
  \item Applying the functor $\FF$ to $\const_z$ produces a morphism
  $g_Z(z) = \FF(\const_z) : \FF(T) \to \FF(Z)$ in $\mathcal{C}$. Viewed as an element of the hom-set,
  $g_Z(z) \in \Hom_{\mathcal{C}}(\FF(T), \FF(Z)) = \CYF(\FF(T))(\FF(Z))$.
\vspace{4 pt}
  \item Consequently, the assignment $z \mapsto g_Z(z)$ is a well-defined \emph{function in $\Set_\ell$}
  from the set $Z \in \mathrm{ob}(\Set_\ell)$ to the hom-set
  $\Hom_{\mathcal{C}}(\FF(T), \FF(Z)) \in \mathrm{ob}(\Set_\ell)$:
  \[
    g_Z \;\in\; \Hom_{\Set_\ell}\Bigl(Z,\, \Hom_{\mathcal{C}}\bigl(\FF(T), \FF(Z)\bigr)\Bigr)
    \;=\; \Hom_{\Set_\ell}\Bigl(Z,\, \CYF(\FF(T))\bigl(\FF(Z)\bigr)\Bigr).
  \]
\vspace{4 pt}
  \item Finally, because $g_Z$ is a morphism in $\Set_\ell$, we can apply the functor
  $\FF : \Set_\ell \to \mathcal{C}$ a \emph{second} time, now to the function $g_Z$ itself, to obtain
  a morphism in $\mathcal{C}$:
  \[
    \FF(g_Z) \;:\; \FF(Z) \longrightarrow \FF\Bigl(\Hom_{\mathcal{C}}\bigl(\FF(T), \FF(Z)\bigr)\Bigr)
    \;=\; G\bigl(\FF(Z)\bigr)
    \quad\text{in }\mathcal{C}.
  \]
\end{enumerate}
Notice that $\FF(g_Z)$ has the exact same domain $\FF(Z)$ and codomain $G(\FF(Z))$ in $\mathcal{C}$ as
the component $\eta_{\FF(Z)} : \FF(Z) \to G(\FF(Z))$ of any natural transformation
$\eta : \Id_{\mathcal{C}} \Rightarrow G$. Equation~\eqref{eq:eta-eq} below simply asserts that these two
morphisms in $\Hom_{\mathcal{C}}(\FF(Z), G(\FF(Z)))$ are equal.
\end{remark}

\begin{lemma}\label{lem:g-id}
For every element $t \in T$, $g_T(t) = \id_{\FF(T)}$. In particular, $g_T(\ttstar) = \id_{\FF(T)}$.
\end{lemma}

\begin{proof}
Since $T = \{\ttstar\}$ has a unique inhabitant, any element $t \in T$ satisfies $t = \ttstar$, and
the constant function $\const_t = (\_ \mapsto t) : T \to T$ is extensionally equal to the identity
function $(z \mapsto z) = \id_T : T \to T$. By functoriality of $\FF$ (preservation of identities),
\[
  g_T(t) \;=\; \FF(\const_t) \;=\; \FF(\id_T) \;=\; \id_{\FF(T)}.
\]
Instantiating at $t = \ttstar$ gives $g_T(\ttstar) = \id_{\FF(T)}$.
\end{proof}

\begin{definition}\label{def:bfc}
A \emph{Basic Functional Category} over $(\mathcal{C}, \FF)$ consists of:
\begin{enumerate}[label=(\roman*),nosep]
  \item A ``monad unit-like'' natural transformation $\eta : \Id_{\mathcal{C}} \Rightarrow G$ between
  the identity endofunctor $\Id_{\mathcal{C}} : \mathcal{C} \to \mathcal{C}$ and the global-like
  endofunctor $G : \mathcal{C} \to \mathcal{C}$, with component morphisms
  \[
    \eta_Y \;:\; Y \longrightarrow G(Y)
    \qquad\bigl(\text{i.e., }\eta_Y \in \Hom_{\mathcal{C}}(Y, G(Y))\bigr)
  \]
  for each $Y \in \mathrm{ob}(\mathcal{C})$, satisfying the naturality condition
  $f \star \eta_X = \eta_Y \star G(f)$ for every morphism $f : Y \to X$ in $\mathcal{C}$;
\vspace{4 pt}  
  \item The general parametric equational law: for every set $Z \in \mathrm{ob}(\Set_\ell)$,
  \begin{equation}\label{eq:eta-eq}
    \eta_{\FF(Z)} \;=\; \FF(g_Z)
    \quad:\quad
    \FF(Z) \longrightarrow G\bigl(\FF(Z)\bigr)
    \quad\text{in }\mathcal{C},
  \end{equation}
  where $g_Z : Z \to \Hom_{\mathcal{C}}(\FF(T), \FF(Z))$ is the global-element generator function in
  $\Set_\ell$ from Definition~\ref{def:cyef}(v) and Remark~\ref{rem:g-typing}.
\end{enumerate}
\end{definition}

\begin{proposition}\label{prop:eta-eq-T}
In any Basic Functional Category $(\mathcal{C}, \FF)$, let $H_T \coloneqq \Hom_{\mathcal{C}}(\FF(T), \FF(T))
\in \mathrm{ob}(\Set_\ell)$. Then the component of $\eta$ at $\FF(T)$ satisfies the terminal unit law:
\begin{equation}\label{eq:eta-eq-T}
  \eta_{\FF(T)} \;=\; g_{H_T}\bigl(g_T(\ttstar)\bigr) \;=\; g_{H_T}\bigl(\id_{\FF(T)}\bigr)
  \quad:\quad
  \FF(T) \longrightarrow G\bigl(\FF(T)\bigr),
\end{equation}
or in subscript-free shorthand, $\eta_{\FF(T)} = g(g(\ttstar)) = g(\id_{\FF(T)})$.
\end{proposition}

\begin{proof}
Instantiating the general parametric law~\eqref{eq:eta-eq} at $Z = T$ gives
$\eta_{\FF(T)} = \FF(g_T)$, where $g_T : T \to H_T$ with $H_T = \Hom_{\mathcal{C}}(\FF(T), \FF(T))$.
Because $T = \{\ttstar\}$ is the terminal set, every $t \in T$ equals $\ttstar$, so $g_T : T \to H_T$
is equal to the constant function $\const_{g_T(\ttstar)} = (\_ \mapsto g_T(\ttstar)) : T \to H_T$. By
Definition~\ref{def:cyef}(v) applied to the set $H_T \in \mathrm{ob}(\Set_\ell)$ and the element
$g_T(\ttstar) \in H_T$, we have
\[
  \eta_{\FF(T)} \;=\; \FF(g_T) \;=\; \FF\bigl(\const_{g_T(\ttstar)}\bigr) \;=\; g_{H_T}\bigl(g_T(\ttstar)\bigr).
\]
Finally, $g_T(\ttstar) = \id_{\FF(T)}$ by Lemma~\ref{lem:g-id}, so
$g_{H_T}(g_T(\ttstar)) = g_{H_T}(\id_{\FF(T)})$.
\end{proof}

\section{Functional Categories and the CoEndoYoneda Isomorphism}\label{sec:coendoyoneda}

\begin{definition}\label{def:fc}
A \emph{Functional Category} over $(\mathcal{C}, \FF)$ is a Basic Functional Category over
$(\mathcal{C}, \FF)$ equipped with:
\begin{enumerate}[label=(\roman*),nosep]
  \item A ``monad multiplication-like'' natural transformation $\mu : G \circ G \Rightarrow G$, with
  component morphisms
  \[
    \mu_Y \;:\; G\bigl(G(Y)\bigr) \longrightarrow G(Y)
    \qquad\bigl(\text{i.e., }\mu_Y \in \Hom_{\mathcal{C}}(G(G(Y)), G(Y))\bigr)
  \]
  for each $Y \in \mathrm{ob}(\mathcal{C})$, satisfying the naturality condition
  $G(G(f)) \star \mu_X = \mu_Y \star G(f)$ for every morphism $f : Y \to X$ in $\mathcal{C}$;
\vspace{4 pt}
  \item T4 pthe monadic left identity law: for every object $Z \in \mathrm{ob}(\mathcal{C})$,
  \begin{equation}\label{eq:id-left}
    \eta_{G(Z)} \star \mu_Z \;=\; \id_{G(Z)}
    \quad:\quad
    G(Z) \longrightarrow G(Z).
  \end{equation}
\end{enumerate}
\end{definition}

\begin{remark}
Notice that neither the right unit law $G(\eta_Z) \star \mu_Z = \id_{G(Z)}$ nor the associativity law
$\mu_{G(Z)} \star \mu_Z = G(\mu_Z) \star \mu_Z$ is required to establish the CoEndoYoneda isomorphism
below: only the left identity law~\eqref{eq:id-left} is needed.
\end{remark}

\subsection{Statement and Construction of the CoEndoYoneda Isomorphism}

Fix a Functional Category $(\mathcal{C}, \FF)$, an object $Z \in \mathrm{ob}(\mathcal{C})$, and an
endofunctor $\EF : \mathcal{C} \to \mathcal{C}$. Set:
\[
  G\circ\EF \;\coloneqq\; G \circ \EF : \mathcal{C} \to \mathcal{C},
  \qquad
  \CYF\circ G\circ\EF \;\coloneqq\; \CYF(\FF(T)) \circ G \circ \EF : \mathcal{C} \to \Set_\ell.
\]
Observe that the set $(\CYF\circ G\circ\EF)(Z)$ is the hom-set of global-like elements of
$(G \circ \EF)(Z)$:
\[
  (\CYF\circ G\circ\EF)(Z) \;=\; \Hom_{\mathcal{C}}\bigl(\FF(T),\, G(\EF(Z))\bigr).
\]
In the classical CoYoneda lemma~\eqref{eq:classical-coyoneda}, the forward isomorphism evaluates a
natural transformation $\tau : \CYF(Z) \Rightarrow F$ on the identity morphism
$\id_Z \in \Hom_{\mathcal{C}}(Z, Z) = \CYF(Z)(Z)$. In our pointfree setting, we lift the element
$\id_Z \in \Hom_{\mathcal{C}}(Z, Z)$ to a global-element morphism $\iotaZ : \FF(T) \to \CYEF(Z)(Z)$ in
$\mathcal{C}$ via the global-element generator~\eqref{eq:def-g} on the set
$\Hom_{\mathcal{C}}(Z, Z) \in \mathrm{ob}(\Set_\ell)$:
\begin{equation}\label{eq:def-iotaZ}
  \iotaZ \;\coloneqq\; g_{\Hom_{\mathcal{C}}(Z, Z)}(\id_Z) \;=\; \FF(\const_{\id_Z})
  \;:\; \FF(T) \longrightarrow \CYEF(Z)(Z),
\end{equation}
where $\const_{\id_Z} = (\_ \mapsto \id_Z) : T \to \Hom_{\mathcal{C}}(Z, Z)$ and
$\CYEF(Z)(Z) = \FF(\Hom_{\mathcal{C}}(Z, Z))$.

\begin{definition}\label{def:witness-maps}
We define two maps between $\NatTrans(\CYEF(Z), G \circ \EF)$ and $(\CYF\circ G\circ\EF)(Z)$:
\begin{enumerate}[label=(\roman*),nosep]
  \item \textbf{Forward extraction map} $\mathrm{to} = (\tau_Z \mapsto f_Z)$: given a natural
  transformation $\tau_Z : \CYEF(Z) \Rightarrow G \circ \EF$, define the morphism
  \begin{equation}\label{eq:tau-to-f}
    \tau_Z \to f_Z \;\coloneqq\; \iotaZ \star (\tau_Z)_Z
    \;:\; \FF(T) \longrightarrow G\bigl(\EF(Z)\bigr)
    \quad\text{in }\Hom_{\mathcal{C}}\bigl(\FF(T),\, G(\EF(Z))\bigr).
  \end{equation}
  \item \textbf{Backward reconstruction map} $\mathrm{from} = (f_Z \mapsto \tau_Z)$: given
  $f_Z \in \Hom_{\mathcal{C}}(\FF(T), G(\EF(Z)))$, define for each object $Y \in \mathrm{ob}(\mathcal{C})$
  the auxiliary function in $\Set_\ell$
  \[
    \alpha_{f_Z, Y} : \Hom_{\mathcal{C}}(Z, Y) \longrightarrow \Hom_{\mathcal{C}}\bigl(\FF(T),\, G(\EF(Y))\bigr),
    \qquad
    z2y \longmapsto f_Z \star (G \circ \EF)(z2y),
  \]
  and set the component morphism $(f_Z \to \tau_Z)_Y : \CYEF(Z)(Y) \to G(\EF(Y))$ in $\mathcal{C}$ to be
  \begin{equation}\label{eq:f-to-tau}
    (f_Z \to \tau_Z)_Y \;\coloneqq\; \FF(\alpha_{f_Z, Y}) \star \mu_{\EF(Y)}
    \;:\; \CYEF(Z)(Y) \longrightarrow G\bigl(\EF(Y)\bigr).
  \end{equation}
\end{enumerate}
\end{definition}

\begin{theorem}\label{thm:coendoyoneda}
In any Functional Category $(\mathcal{C}, \FF)$, for every object $Z \in \mathrm{ob}(\mathcal{C})$ and
every endofunctor $\EF : \mathcal{C} \to \mathcal{C}$:
\begin{enumerate}[label=(\roman*),nosep]
  \item For every $f_Z \in (\CYF\circ G\circ\EF)(Z)$, the family $(f_Z \to \tau_Z)_Y$ defined in
  \eqref{eq:f-to-tau} is natural in $Y$, defining a natural transformation
  $(f_Z \to \tau_Z) : \CYEF(Z) \Rightarrow G \circ \EF$.
\vspace{4 pt}
  \item \textbf{Section:} For every $f_Z \in (\CYF\circ G\circ\EF)(Z)$,
  \[
    (f_Z \to \tau_Z) \to f_Z \;=\; f_Z.
  \]
\vspace{4 pt}
  \item \textbf{Retraction:} For every
  $\tau_Z \in \NatTrans(\CYEF(Z), G \circ \EF)$,
  \[
    (\tau_Z \to f_Z) \to \tau_Z \;=\; \tau_Z.
  \]
\end{enumerate}
Consequently, $(\tau_Z \to f_Z)$ and $(f_Z \to \tau_Z)$ form a bijection (isomorphism of types)
\[
  \NatTrans\bigl(\CYEF(Z),\, G \circ \EF\bigr)
  \;\cong\;
  \Hom_{\mathcal{C}}\bigl(\FF(T),\, (G \circ \EF)(Z)\bigr).
\]
\end{theorem}

\begin{proof}
We follow the equational chains formalized in~\cite{CubicalAgdaRepo} and mechanized
in~\cite{PalomarLeanRepo}.

\textbf{Part (i): Naturality of $f_Z \to \tau_Z$.}
First, for any morphisms $y2x : Y \to X$ and $z2y : Z \to Y$ in $\mathcal{C}$, functoriality of
$G \circ \EF$ and associativity of composition in $\mathcal{C}$ give:
\begin{align*}
  f_Z \star (G\circ\EF)(z2y \star y2x)
  &= f_Z \star \bigl((G\circ\EF)(z2y) \star (G\circ\EF)(y2x)\bigr) \\
  &= \bigl(f_Z \star (G\circ\EF)(z2y)\bigr) \star (G\circ\EF)(y2x).
\end{align*}
Now let $y2x : Y \to X$ in $\mathcal{C}$. Using $\CYEF(Z)(y2x) = \FF(z2y \mapsto z2y \star y2x)$,
functoriality of $\FF$ on composite functions, the definition
$(G \circ G)(\EF(y2x)) = \FF(h \mapsto h \star (G\circ\EF)(y2x))$, and naturality of $\mu$ at
$\EF(y2x)$:
\begin{align*}
  &\CYEF(Z)(y2x) \star (f_Z \to \tau_Z)_X \\
  &\quad= \CYEF(Z)(y2x) \star \bigl(\FF(\alpha_{f_Z, X}) \star \mu_{\EF(X)}\bigr)
   \;=\; \bigl(\CYEF(Z)(y2x) \star \FF(\alpha_{f_Z, X})\bigr) \star \mu_{\EF(X)} \\
  &\quad= \FF\bigl(z2y \mapsto f_Z \star (G\circ\EF)(z2y \star y2x)\bigr) \star \mu_{\EF(X)} \\
  &\quad= \FF\bigl(z2y \mapsto (f_Z \star (G\circ\EF)(z2y)) \star (G\circ\EF)(y2x)\bigr) \star \mu_{\EF(X)} \\
  &\quad= \bigl(\FF(\alpha_{f_Z, Y}) \star (G \circ G)(\EF(y2x))\bigr) \star \mu_{\EF(X)}
   \;=\; \FF(\alpha_{f_Z, Y}) \star \bigl((G \circ G)(\EF(y2x)) \star \mu_{\EF(X)}\bigr) \\
  &\quad= \FF(\alpha_{f_Z, Y}) \star \bigl(\mu_{\EF(Y)} \star (G\circ\EF)(y2x)\bigr) \\
  &\quad= \bigl(\FF(\alpha_{f_Z, Y}) \star \mu_{\EF(Y)}\bigr) \star (G\circ\EF)(y2x)
   \;=\; (f_Z \to \tau_Z)_Y \star (G\circ\EF)(y2x).
\end{align*}

\textbf{Part (ii): Section.}
Let $f_Z : \FF(T) \to G(\EF(Z))$. Using $(G\circ\EF)(\id_Z) = \id_{G(\EF(Z))}$, Lemma~\ref{lem:g-id}
($\id_{\FF(T)} = g_T(t)$ for $t \in T$), functoriality of $\FF$ on $g_T$ followed by
$(h \mapsto h \star f_Z)$, the parametric unit law~\eqref{eq:eta-eq} at $Z = T$
($\FF(g_T) = \eta_{\FF(T)}$), naturality of $\eta$ on $f_Z$, and the left identity
law~\eqref{eq:id-left}:
\begin{align*}
  ((f_Z \to \tau_Z) \to f_Z)
  &= \iotaZ \star \bigl(\FF(z2z \mapsto f_Z \star (G\circ\EF)(z2z)) \star \mu_{\EF(Z)}\bigr) \\
  &= \bigl(g(\id_Z) \star \FF(z2z \mapsto f_Z \star (G\circ\EF)(z2z))\bigr) \star \mu_{\EF(Z)} \\
  &= g\bigl(f_Z \star (G\circ\EF)(\id_Z)\bigr) \star \mu_{\EF(Z)} \\
  &= g(f_Z \star \id_{G(\EF(Z))}) \star \mu_{\EF(Z)} \;=\; g(f_Z) \star \mu_{\EF(Z)}
   \;=\; g(\id_{\FF(T)} \star f_Z) \star \mu_{\EF(Z)} \\
  &= \FF\bigl(t \mapsto g_T(t) \star f_Z\bigr) \star \mu_{\EF(Z)} \\
  &= \bigl(\FF(g_T) \star G(f_Z)\bigr) \star \mu_{\EF(Z)} \\
  &= \FF(g_T) \star \bigl(G(f_Z) \star \mu_{\EF(Z)}\bigr) \\
  &= \eta_{\FF(T)} \star \bigl(G(f_Z) \star \mu_{\EF(Z)}\bigr) \\
  &= \bigl(\eta_{\FF(T)} \star G(f_Z)\bigr) \star \mu_{\EF(Z)} \\
  &= \bigl(f_Z \star \eta_{G(\EF(Z))}\bigr) \star \mu_{\EF(Z)} \\
  &= f_Z \star \bigl(\eta_{G(\EF(Z))} \star \mu_{\EF(Z)}\bigr)
   \;=\; f_Z \star \id_{G(\EF(Z))} \;=\; f_Z.
\end{align*}

\textbf{Part (iii): Retraction.}
Let $\tau_Z : \CYEF(Z) \Rightarrow G \circ \EF$ and $Y \in \mathrm{ob}(\mathcal{C})$. Applying
associativity, naturality of $\tau_Z$ along $z2y : Z \to Y$, the identity
$\iotaZ \star \CYEF(Z)(z2y) = g(\id_Z \star z2y) = g(z2y)$, the parametric unit
law~\eqref{eq:eta-eq} at the hom-set $H_{Z,Y} \coloneqq \Hom_{\mathcal{C}}(Z, Y)$
($\FF(g_{H_{Z,Y}}) = \eta_{\FF(\Hom_{\mathcal{C}}(Z, Y))} = \eta_{\CYEF(Z)(Y)}$), naturality of $\eta$
along $(\tau_Z)_Y$, and the left identity law~\eqref{eq:id-left}:
\begin{align*}
  ((\tau_Z \to f_Z) \to \tau_Z)_Y
  &= \FF\bigl(z2y \mapsto (\iotaZ \star (\tau_Z)_Z) \star (G\circ\EF)(z2y)\bigr) \star \mu_{\EF(Y)} \\
  &= \FF\bigl(z2y \mapsto \iotaZ \star ((\tau_Z)_Z \star (G\circ\EF)(z2y))\bigr) \star \mu_{\EF(Y)} \\
  &= \FF\bigl(z2y \mapsto \iotaZ \star (\CYEF(Z)(z2y) \star (\tau_Z)_Y)\bigr) \star \mu_{\EF(Y)} \\
  &= \FF\bigl(z2y \mapsto (\iotaZ \star \CYEF(Z)(z2y)) \star (\tau_Z)_Y\bigr) \star \mu_{\EF(Y)} \\
  &= \FF\bigl(z2y \mapsto g(\id_Z \star z2y) \star (\tau_Z)_Y\bigr) \star \mu_{\EF(Y)} \\
  &= \FF\bigl(z2y \mapsto g(z2y) \star (\tau_Z)_Y\bigr) \star \mu_{\EF(Y)} \\
  &= \bigl(\FF(g_{H_{Z,Y}}) \star G((\tau_Z)_Y)\bigr) \star \mu_{\EF(Y)} \\
  &= \FF(g_{H_{Z,Y}}) \star \bigl(G((\tau_Z)_Y) \star \mu_{\EF(Y)}\bigr) \\
  &= \eta_{\CYEF(Z)(Y)} \star \bigl(G((\tau_Z)_Y) \star \mu_{\EF(Y)}\bigr) \\
  &= \bigl(\eta_{\CYEF(Z)(Y)} \star G((\tau_Z)_Y)\bigr) \star \mu_{\EF(Y)} \\
  &= \bigl((\tau_Z)_Y \star \eta_{G(\EF(Y))}\bigr) \star \mu_{\EF(Y)} \\
  &= (\tau_Z)_Y \star \bigl(\eta_{G(\EF(Y))} \star \mu_{\EF(Y)}\bigr)
   \;=\; (\tau_Z)_Y \star \id_{G(\EF(Y))} \;=\; (\tau_Z)_Y.
\end{align*}
By extensionality of natural transformations, $((\tau_Z \to f_Z) \to \tau_Z) = \tau_Z$.
\end{proof}

\subsection{Connecting the CoEndoYoneda Isomorphism to $\eta$ at $\FF(T)$}

When we specialize Theorem~\ref{thm:coendoyoneda} to $\EF = \Id_{\mathcal{C}}$ and $Z = \FF(T)$, the
domain functor of the left-hand side becomes $\CYEF(\FF(T)) = G$, and the codomain functor becomes
$G \circ \Id_{\mathcal{C}} = G$. Hence the left-hand side contains the canonical identity natural
transformation $\id_G : G \Rightarrow G$.

\begin{proposition}\label{prop:origin-eta}
Define the morphism
\[
  \eta_{\FF(T)}^{\mathrm{CY}} \;\coloneqq\; (\id_G \to f_{\FF(T)})
  \;:\; \FF(T) \longrightarrow G\bigl(\FF(T)\bigr)
\]
as the image of $\id_G$ under the forward map~\eqref{eq:tau-to-f} of the CoEndoYoneda isomorphism at
$\EF = \Id_{\mathcal{C}}$ and $Z = \FF(T)$. Then in any Basic Functional Category,
\[
  \eta_{\FF(T)}^{\mathrm{CY}} \;=\; \eta_{\FF(T)}.
\]
\end{proposition}

\begin{proof}
Unfolding \eqref{eq:tau-to-f} at $\tau = \id_G$, and applying right identity, Lemma~\ref{lem:g-id},
and Proposition~\ref{prop:eta-eq-T}:
\[
  \eta_{\FF(T)}^{\mathrm{CY}}
  \;=\; \iota_{\FF(T)} \star \id_{G(\FF(T))}
  \;=\; g(\id_{\FF(T)})
  \;=\; g\bigl(g(\ttstar)\bigr)
  \;=\; \eta_{\FF(T)}.
  \qedhere
\]
\end{proof}

\subsection{Canonical Concrete Instance: $\Set$ as a Functional Category}

\begin{proposition}\label{prop:set-fc}
In the category $\mathcal{C} = \Set_\ell$ with $\FF = \Id_{\Set_\ell}$, we have
$G(Y) = \CYEF(T)(Y) = (T \to Y) \cong Y$. Define for each $Z \in \mathrm{ob}(\Set_\ell)$:
\[
  \eta_Z(z) \;\coloneqq\; \const_z \;:\; T \to Z,
  \qquad
  \mu_Z(h) \;\coloneqq\; (t \mapsto h(t)(t)) \;:\; T \to Z
  \quad\text{for } h : T \to (T \to Z).
\]
Then $\eta : \Id_{\Set_\ell} \Rightarrow G$ and $\mu : G \circ G \Rightarrow G$ are natural
transformations, and the laws \eqref{eq:eta-eq}, \eqref{eq:eta-eq-T}, and \eqref{eq:id-left} all hold
definitionally.
\end{proposition}

\begin{proof}
Since $\FF = \Id_{\Set_\ell}$, $\FF(g_Z)(z) = g_Z(z) = \const_z = \eta_Z(z)$ definitionally, proving
\eqref{eq:eta-eq} (and hence \eqref{eq:eta-eq-T}) by reflexivity. For any $h : T \to Z$ and $t \in T$,
$(\eta_{G(Z)} \star \mu_Z)(h)(t) = \mu_Z(\_ \mapsto h)(t) = (\_ \mapsto h)(t)(t) = h(t)$, so
$\eta_{G(Z)} \star \mu_Z = \id_{G(Z)}$ holds definitionally by $\beta\eta$-reduction.
\end{proof}

\subsection{Relationship to the Classical CoYoneda Lemma: Naturality, Universality, Generalization, and Specialization}\label{subsec:coyoneda-comparison}

It is illuminating to compare the CoEndoYoneda Isomorphism Theorem (Theorem~\ref{thm:coendoyoneda})
with the classical CoYoneda lemma~\eqref{eq:classical-coyoneda} across four foundational dimensions:

\begin{enumerate}[label=(\arabic*),nosep]
  \item \textbf{Endofunctor coincidence on $\mathcal{C} = \Set$ and generalization to $(\mathcal{C}, \FF)$:}
  In the classical CoYoneda lemma $\NatTrans(\CYF(Z), F) \cong F(Z)$, the domain functor
  $\CYF(Z) = \Hom_{\mathcal{C}}(Z, -) : \mathcal{C} \to \Set$ is always $\Set$-valued. In the special
  case where the base category $\mathcal{C}$ is $\Set$ itself, $\CYF(Z) : \Set \to \Set$ and
  $F = \EF : \Set \to \Set$ are both \emph{endofunctors} on $\Set$. When we instantiate our
  CoEndoYoneda Isomorphism Theorem~\ref{thm:coendoyoneda} at the canonical Functional Category
  $(\mathcal{C}, \FF) = (\Set, \Id_{\Set})$ of Proposition~\ref{prop:set-fc}, we have
  $\CYEF(Z) = \Id_{\Set} \circ \CYF(Z) = \CYF(Z)$, $G = (T \to -) \cong \Id_{\Set}$, and
  \[
    \NatTrans\bigl(\CYEF(Z),\, G \circ \EF\bigr)
    \;\cong\;
    \Hom_{\Set}\bigl(T,\, G(\EF(Z))\bigr)
    \;=\;
    \bigl(T \to (T \to \EF(Z))\bigr)
    \;\cong\;
    \EF(Z),
  \]
  which recovers the classical CoYoneda lemma for endofunctors on $\Set$. Conversely, for an arbitrary
  Functional Category $(\mathcal{C}, \FF)$, Theorem~\ref{thm:coendoyoneda} \emph{generalizes} this
  endofunctor CoYoneda isomorphism from $\Set$ to endofunctors on $\mathcal{C}$ by lifting $\CYF(Z)$ to
  $\CYEF(Z) = \FF \circ \CYF(Z) : \mathcal{C} \to \mathcal{C}$ and replacing set-theoretic elements of
  $F(Z)$ by global-like morphisms $\FF(T) \to G(\EF(Z))$ in $\mathcal{C}$.

  \vspace{3pt}
  \item \textbf{Pointfree universality and the universal morphism $\iotaZ$:}
  In the classical CoYoneda lemma, the pair $(Z, \id_Z \in \CYF(Z)(Z))$ is the universal element of the
  representable presheaf $\CYF(Z)$: every natural transformation $\tau : \CYF(Z) \Rightarrow F$ is
  uniquely determined by the single element $x = \tau_Z(\id_Z) \in F(Z)$. In a general category
  $\mathcal{C}$, however, the object $\CYEF(Z)(Z) = \FF(\Hom_{\mathcal{C}}(Z, Z)) \in \mathrm{ob}(\mathcal{C})$
  is an abstract object of $\mathcal{C}$ rather than a set, so it has no ``elements'' in the pointful
  sense. In Theorem~\ref{thm:coendoyoneda}, the role of the universal element $\id_Z$ is played
  pointfreely by the universal global-element morphism
  \[
    \iotaZ \;=\; g_{\Hom_{\mathcal{C}}(Z, Z)}(\id_Z) \;=\; \FF(\const_{\id_Z})
    \;:\; \FF(T) \longrightarrow \CYEF(Z)(Z)
    \quad\text{in }\mathcal{C},
  \]
  and evaluation at $\id_Z$ is replaced by diagrammatic precomposition $\tau_Z \mapsto \iotaZ \star (\tau_Z)_Z$.
  Moreover, Proposition~\ref{prop:origin-eta} shows that the universal element corresponding to the
  identity natural transformation $\id_G : G \Rightarrow G$ at $Z = \FF(T)$ is precisely the terminal
  unit component $\eta_{\FF(T)} = \iota_{\FF(T)} \star \id_{G(\FF(T))} = g(\id_{\FF(T)})$.

  \vspace{3pt}
  \item \textbf{Role of the Functional Category structure $(\eta, \mu)$ in reconstruction:}
  In the classical CoYoneda lemma, reconstructing $\tau_Y : \Hom_{\mathcal{C}}(Z, Y) \to F(Y)$ from
  $x \in F(Z)$ sends $z2y \mapsto F(z2y)(x)$. In our setting, given $f_Z : \FF(T) \to G(\EF(Z))$,
  postcomposing with $(G \circ \EF)(z2y)$ yields a function
  $\alpha_{f_Z, Y} : \Hom_{\mathcal{C}}(Z, Y) \to \Hom_{\mathcal{C}}(\FF(T), G(\EF(Y)))$ in $\Set$.
  Applying $\FF$ produces a morphism
  \[
    \FF(\alpha_{f_Z, Y}) \;:\; \CYEF(Z)(Y) \longrightarrow G\bigl(G(\EF(Y))\bigr)
    \quad\text{in }\mathcal{C},
  \]
  which lands in $G(G(\EF(Y)))$ with \emph{two} layers of $G$. The multiplication-like transformation
  $\mu_{\EF(Y)} : G(G(\EF(Y))) \to G(\EF(Y))$ flattens these two layers into one, while the parametric
  unit law $\eta_{\FF(Z)} = \FF(g_Z)$ and the left unit law $\eta_{G(X)} \star \mu_X = \id_{G(X)}$ are
  the exact coherence axioms needed to guarantee that flattening after lifting is a two-sided inverse to
  precomposition with $\iotaZ$.

  \vspace{3pt}
  \item \textbf{Comparison with the adjunction perspective:}
  When $\FF : \Set \to \mathcal{C}$ is a left adjoint to the global-elements functor
  $\CYF(\FF(T)) = \Hom_{\mathcal{C}}(\FF(T), -) : \mathcal{C} \to \Set$, natural transformations
  $\FF \circ \CYF(Z) \Rightarrow G \circ \EF$ transpose across the adjunction to natural transformations
  $\CYF(Z) \Rightarrow \CYF(\FF(T)) \circ G \circ \EF$ in $\Set$, where the classical CoYoneda lemma
  applies. Crucially, however, a Functional Category $(\mathcal{C}, \FF)$ \emph{does not assume an
  adjunction} $\FF \dashv \CYF(\FF(T))$ (nor even that $G$ is a full monad, since neither the right unit
  law nor associativity of $\mu$ is required). Instead, Theorem~\ref{thm:coendoyoneda} isolates the
  minimal pointfree equational core---Equations~\eqref{eq:eta-eq} and~\eqref{eq:id-left}---under which
  the CoEndoYoneda bijection holds directly inside $\mathcal{C}$.
\end{enumerate}

\section{Purity, Idempotent Monads, and Kleisli Categories}\label{sec:purity}

We now turn to the connection between Basic Functional Categories and computational purity for monads
on $\Set_\ell$.

\subsection{Monad Formulations and Characterizations of Purity}

In functional programming, a monad $(M, \mathrm{mPure}, \mathrm{mApply})$ on $\Set_\ell$ is given by an
object mapping $M : \mathrm{ob}(\Set_\ell) \to \mathrm{ob}(\Set_\ell)$, a unit family
$\mathrm{mPure}_Z : Z \to M(Z)$, and a Kleisli extension (bind) operation
$\mathrm{mApply} : (Z \to M(Y)) \to M(Z) \to M(Y)$ satisfying the three monad laws:
\begin{align}
  \mathrm{mApply}\;f\;(\mathrm{mPure}_Z\;z) &= f(z)
    &&\text{(left unit law)}, \label{eq:monad-left-unit} \\
  \mathrm{mApply}\;\mathrm{mPure}_Z\;mz &= mz
    &&\text{(right unit law)}, \label{eq:monad-right-unit} \\
  \mathrm{mApply}\;h\;(\mathrm{mApply}\;f\;mz) &= \mathrm{mApply}\;\bigl(z \mapsto \mathrm{mApply}\;h\;(f(z))\bigr)\;mz
    &&\text{(associativity law)}. \label{eq:monad-assoc}
\end{align}
From these, one defines the functorial map operation $\mathrm{mMap}$ (for $f : Z \to Y$) and monadic
multiplication $\mathrm{mJoin}$:
\begin{equation}\label{eq:mmap-mjoin}
  \mathrm{mMap}\;f\;mz \;\coloneqq\; \mathrm{mApply}\;(\mathrm{mPure}_Y \circ f)\;mz,
  \qquad
  \mathrm{mJoin}\;mmz \;\coloneqq\; \mathrm{mApply}\;\id_{M(Z)}\;mmz,
\end{equation}
which satisfy all standard categorical monad and naturality laws, with
$\mathrm{mApply}\;f\;mz = \mathrm{mJoin}(\mathrm{mMap}\;f\;mz)$.

\begin{definition}\label{def:idempotent}
Let $(M, \mathrm{mPure}, \mathrm{mApply})$ be a monad on $\Set_\ell$.
\begin{enumerate}[label=(\roman*),nosep]
  \item $M$ is \emph{idempotent} ($\mathrm{Idempotent}(M)$, cfr.\ $M\eta = \eta M$ in~\cite{Borceux1994,nLabIdempotentMonad})
  if for every $Z \in \mathrm{ob}(\Set_\ell)$ and $mz \in M(Z)$,
  \begin{equation}\label{eq:idempotent}
    \mathrm{mMap}\;\mathrm{mPure}_Z\;mz \;=\; \mathrm{mPure}_{M(Z)}\;mz.
  \end{equation}
  Intuitively, performing an effect twice is the same as performing it once.
  \item $M$ \emph{enforces purity} ($\mathrm{IsPure}(M)$) if for every $Z \in \mathrm{ob}(\Set_\ell)$ and
  $mz \in M(Z)$,
  \begin{equation}\label{eq:ispure}
    \mathrm{mApply}\;(\mathrm{mPure}_{M(Z)} \circ \mathrm{mPure}_Z)\;mz \;=\; \mathrm{mPure}_{M(Z)}\;mz.
  \end{equation}
\end{enumerate}
\end{definition}

\begin{lemma}\label{lem:idem-ispure}
For any monad $M$ on $\Set_\ell$, $\mathrm{Idempotent}(M) \iff \mathrm{IsPure}(M)$.
\end{lemma}

\begin{proof}
Immediate because $\mathrm{mMap}\;\mathrm{mPure}_Z\;mz =
\mathrm{mApply}\;(\mathrm{mPure}_{M(Z)} \circ \mathrm{mPure}_Z)\;mz$ by definition~\eqref{eq:mmap-mjoin}
of $\mathrm{mMap}$.
\end{proof}

\subsection{Equivalence Between Kleisli Basic Functional Categories and Monad Idempotency}

Given a monad $M$ on $\Set_\ell$, its \emph{Kleisli category} $\K(M)$ has:
\begin{itemize}[nosep]
  \item \textbf{Objects:} $\mathrm{ob}(\K(M)) \coloneqq \mathrm{ob}(\Set_\ell)$;
\vspace{4 pt}  
  \item \textbf{Morphisms:} $\Hom_{\K(M)}(Z, Y) \coloneqq (Z \to M(Y))$;
\vspace{4 pt}  
  \item \textbf{Identities:} $\id_Z^{\K(M)} \coloneqq \mathrm{mPure}_Z : Z \to M(Z)$;
\vspace{4 pt}  
  \item \textbf{Kleisli composition} (denoted $\starK$ to distinguish it from ordinary function
  composition in $\Set_\ell$): for $f \in \Hom_{\K(M)}(Z, Y)$ and $h \in \Hom_{\K(M)}(Y, X)$,
  \begin{equation}\label{eq:kleisli-comp}
    (f \starK h)(z) \;\coloneqq\; \mathrm{mApply}\;h\;\bigl(f(z)\bigr) \;\in\; M(X);
  \end{equation}
\vspace{4 pt}  
  \item \textbf{Canonical functional functor} $\KFF : \Set_\ell \to \K(M)$ (identity on objects):
  \begin{equation}\label{eq:kff}
    \KFF(Z) \;\coloneqq\; Z,
    \qquad
    \KFF(f : Z \to Y) \;\coloneqq\; \bigl(z \mapsto \mathrm{mPure}_Y(f(z))\bigr)
    \;\in\; \Hom_{\K(M)}(Z, Y).
  \end{equation}
\end{itemize}
In $(\K(M), \KFF)$, let us explicitly unpack the global-like endofunctor
$G = \CYEF(\KFF(T)) = \KFF \circ \CYF(T) : \K(M) \to \K(M)$ and the global-element generator $g_Z$:
\begin{itemize}[nosep]
  \item For any object $Z \in \mathrm{ob}(\K(M))$, since $\KFF$ is the identity on objects,
  \[
    G(Z) \;=\; \KFF\bigl(\Hom_{\K(M)}(T, Z)\bigr) \;=\; \Hom_{\K(M)}(T, Z) \;=\; \bigl(T \to M(Z)\bigr).
  \]
\vspace{4 pt}  
  \item For any Kleisli morphism $f : Y \to M(X)$ in $\Hom_{\K(M)}(Y, X)$, the function
  $\CYF(T)(f) : (T \to M(Y)) \to (T \to M(X))$ maps $g_{MY} : T \to M(Y)$ to the Kleisli composite
  $g_{MY} \starK f = \bigl(t \mapsto \mathrm{mApply}\;f\;(g_{MY}(t))\bigr)$. Applying $\KFF$ according to
  \eqref{eq:kff} yields
  \begin{equation}\label{eq:kleisli-G-map}
    G(f)(g_{MY}) \;=\; \mathrm{mPure}_{G(X)}\bigl(t \mapsto \mathrm{mApply}\;f\;(g_{MY}(t))\bigr)
    \;\in\; M\bigl(T \to M(X)\bigr).
  \end{equation}
\vspace{4 pt}  
  \item For any set $Z$ and element $z \in Z$, the global-element generator
  $g_Z(z) = \KFF(\const_z) \in \Hom_{\K(M)}(T, Z)$ is the constant function
  $g_Z(z) = \bigl(\_ \mapsto \mathrm{mPure}_Z(z)\bigr) : T \to M(Z)$. Consequently, lifting
  $g_Z : Z \to G(Z)$ via $\KFF$ yields
  \begin{equation}\label{eq:kleisli-kff-g}
    \KFF(g_Z)(z) \;=\; \mathrm{mPure}_{G(Z)}\bigl(g_Z(z)\bigr)
    \;=\; \mathrm{mPure}_{G(Z)}\bigl(\_ \mapsto \mathrm{mPure}_Z(z)\bigr)
    \;\in\; M\bigl(T \to M(Z)\bigr),
  \end{equation}
  and at $Z = T$ with $z = \ttstar$, the terminal law~\eqref{eq:eta-eq-T} states that
  \begin{equation}\label{eq:kleisli-eta-T}
    \eta_T(\ttstar) \;=\; g_{G(T)}\bigl(g_T(\ttstar)\bigr)(\ttstar)
    \;=\; \mathrm{mPure}_{G(T)}\bigl(\_ \mapsto \mathrm{mPure}_T(\ttstar)\bigr).
  \end{equation}
\end{itemize}

\begin{theorem}\label{thm:kleisli-purity}
Let $(M, \mathrm{mPure}, \mathrm{mApply})$ be a monad on $\Set_\ell$. Then $(\K(M), \KFF)$ is a
\emph{Basic Functional Category} if and only if $M$ is an \emph{Idempotent} (pure) monad:
\[
  \mathrm{BasicFunctionalCategory}\bigl(\K(M),\, \KFF\bigr)
  \;\Longleftrightarrow\;
  \mathrm{Idempotent}(M).
\]
More specifically, there is a striking implication associated with the two directions:
\begin{enumerate}[label=(\roman*),nosep]
  \item \textbf{($\Longrightarrow$):} Only naturality of $\eta$ and the terminal specialization
  \eqref{eq:eta-eq-T} ($\eta_T = g_{G(T)}(g_T(\ttstar))$) are used to deduce that $M$ is globally
  idempotent on all sets $Z$.
\vspace{4 pt}  
  \item \textbf{($\Longleftarrow$):} Assuming $M$ is idempotent, the candidate unit
  $\eta_Y(y) \coloneqq \mathrm{mPure}_{G(Y)}(\_ \mapsto \mathrm{mPure}_Y(y))$ is natural and satisfies
  the full parametric law~\eqref{eq:eta-eq} ($\eta_Z = \KFF(g_Z)$) definitionally by reflexivity.
\end{enumerate}
\end{theorem}

\begin{proof}
We provide the complete, self-contained equational proof of both directions, matching the mechanized
derivations in~\cite{CubicalAgdaRepo,PalomarLeanRepo}.

\textbf{Direction ($\Longrightarrow$): $\mathrm{BasicFunctionalCategory}(\K(M), \KFF) \implies \mathrm{Idempotent}(M)$.}
Suppose $(\K(M), \KFF)$ is a Basic Functional Category with unit natural transformation
$\eta : \Id_{\K(M)} \Rightarrow G$, where for each set $Z$, $\eta_Z : Z \to M(T \to M(Z))$ is a
Kleisli morphism in $\Hom_{\K(M)}(Z, G(Z))$.

\emph{Step 1: Proof of the Evaluation Lemma.}
We first prove that for every set $Z \in \mathrm{ob}(\Set_\ell)$ and every monadic computation
$mz \in M(Z)$,
\begin{equation}\label{eq:eval-lemma}
  \mathrm{mApply}\;\eta_Z\;mz \;=\; \mathrm{mPure}_{G(Z)}(\_ \mapsto mz)
  \;\in\; M\bigl(T \to M(Z)\bigr).
\end{equation}
Define the constant Kleisli morphism $t\!\to\!mz \in \Hom_{\K(M)}(T, Z)$ by
\[
  (t\!\to\!mz)(t) \;\coloneqq\; mz \qquad\text{for all } t \in T.
\]
By naturality of $\eta : \Id_{\K(M)} \Rightarrow G$ applied to the Kleisli morphism
$t\!\to\!mz : T \to Z$ in $\K(M)$, we have the equality of Kleisli morphisms in
$\Hom_{\K(M)}(T, G(Z))$:
\begin{equation}\label{eq:eta-nat-kleisli}
  (t\!\to\!mz) \starK \eta_Z \;=\; \eta_T \starK G(t\!\to\!mz).
\end{equation}
We evaluate both sides of \eqref{eq:eta-nat-kleisli} at the unique element $\ttstar \in T$.
By definition~\eqref{eq:kleisli-comp} of Kleisli composition, the left-hand side of
\eqref{eq:eta-nat-kleisli} evaluated at $\ttstar$ is:
\begin{equation}\label{eq:eval-lhs}
  \bigl((t\!\to\!mz) \starK \eta_Z\bigr)(\ttstar)
  \;=\; \mathrm{mApply}\;\eta_Z\;\bigl((t\!\to\!mz)(\ttstar)\bigr)
  \;=\; \mathrm{mApply}\;\eta_Z\;mz.
\end{equation}
On the right-hand side of \eqref{eq:eta-nat-kleisli}, evaluating at $\ttstar$ and applying the
Kleisli composition formula~\eqref{eq:kleisli-comp}, the terminal unit law~\eqref{eq:kleisli-eta-T},
the left unit monad law~\eqref{eq:monad-left-unit}, the action of $G$ in~\eqref{eq:kleisli-G-map}, and
a second application of the left unit monad law~\eqref{eq:monad-left-unit} yields:
\begin{align}
  &\bigl(\eta_T \starK G(t\!\to\!mz)\bigr)(\ttstar) \nonumber \\
  &\quad= \mathrm{mApply}\;\bigl(G(t\!\to\!mz)\bigr)\;\bigl(\eta_T(\ttstar)\bigr)
    &&\text{(by \eqref{eq:kleisli-comp})} \nonumber \\
  &\quad= \mathrm{mApply}\;\bigl(G(t\!\to\!mz)\bigr)\;\bigl(\mathrm{mPure}_{G(T)}(\_ \mapsto \mathrm{mPure}_T(\ttstar))\bigr)
    &&\text{(by \eqref{eq:kleisli-eta-T})} \nonumber \\
  &\quad= G(t\!\to\!mz)\bigl(\_ \mapsto \mathrm{mPure}_T(\ttstar)\bigr)
    &&\text{(by \eqref{eq:monad-left-unit})} \nonumber \\
  &\quad= \mathrm{mPure}_{G(Z)}\bigl(t \mapsto \mathrm{mApply}\;(t\!\to\!mz)\;((\_ \mapsto \mathrm{mPure}_T(\ttstar))(t))\bigr)
    &&\text{(by \eqref{eq:kleisli-G-map})} \nonumber \\
  &\quad= \mathrm{mPure}_{G(Z)}\bigl(t \mapsto \mathrm{mApply}\;(t\!\to\!mz)\;(\mathrm{mPure}_T(\ttstar))\bigr)
    &&\text{($\beta$-red.)} \nonumber \\
  &\quad= \mathrm{mPure}_{G(Z)}\bigl(t \mapsto (t\!\to\!mz)(\ttstar)\bigr)
    &&\text{(by \eqref{eq:monad-left-unit})} \nonumber \\
  &\quad= \mathrm{mPure}_{G(Z)}(\_ \mapsto mz)
    &&\text{(def.\ $t\!\to\!mz$)}. \label{eq:eval-rhs}
\end{align}
Combining \eqref{eq:eta-nat-kleisli}, \eqref{eq:eval-lhs}, and \eqref{eq:eval-rhs} establishes the
Evaluation Lemma~\eqref{eq:eval-lemma}.

\emph{Step 2: Point Evaluation Homomorphism and Pointwise Idempotency Lemma.}
For each set $Z$, define the evaluation map $\mathrm{mgz\!\to\!mmz} : M(G(Z)) \to M(M(Z))$ by
\begin{equation}\label{eq:mgz-to-mmz}
  \mathrm{mgz\!\to\!mmz}(m_{t\to mz})
  \;\coloneqq\; \mathrm{mApply}\;\bigl(t2mz \mapsto \mathrm{mPure}_{M(Z)}(t2mz(\ttstar))\bigr)\;m_{t\to mz}.
\end{equation}
Note that for any function $h : T \to M(Z)$, the left unit monad law~\eqref{eq:monad-left-unit} gives
\begin{align}
  \mathrm{mgz\!\to\!mmz}\bigl(\mathrm{mPure}_{G(Z)}(h)\bigr)
  &= \mathrm{mApply}\;\bigl(t2mz \mapsto \mathrm{mPure}_{M(Z)}(t2mz(\ttstar))\bigr)\;\bigl(\mathrm{mPure}_{G(Z)}(h)\bigr) \nonumber \\
  &= \mathrm{mPure}_{M(Z)}\bigl(h(\ttstar)\bigr). \label{eq:mgz-on-pure}
\end{align}
Now, for any element $z \in Z$, applying the left unit monad law~\eqref{eq:monad-left-unit} to write
$\eta_Z(z) = \mathrm{mApply}\;\eta_Z\;(\mathrm{mPure}_Z(z))$, then applying the Evaluation
Lemma~\eqref{eq:eval-lemma} to $mz = \mathrm{mPure}_Z(z)$, and finally applying~\eqref{eq:mgz-on-pure}
to $h = (\_ \mapsto \mathrm{mPure}_Z(z))$ gives:
\begin{align}
  \mathrm{mgz\!\to\!mmz}\bigl(\eta_Z(z)\bigr)
  &= \mathrm{mgz\!\to\!mmz}\Bigl(\mathrm{mApply}\;\eta_Z\;\bigl(\mathrm{mPure}_Z(z)\bigr)\Bigr)
    &&\text{(by \eqref{eq:monad-left-unit})} \nonumber \\
  &= \mathrm{mgz\!\to\!mmz}\Bigl(\mathrm{mPure}_{G(Z)}\bigl(\_ \mapsto \mathrm{mPure}_Z(z)\bigr)\Bigr)
    &&\text{(by \eqref{eq:eval-lemma})} \nonumber \\
  &= \mathrm{mPure}_{M(Z)}\bigl(\mathrm{mPure}_Z(z)\bigr)
    &&\text{(by \eqref{eq:mgz-on-pure})}. \label{eq:idempotent-lemma}
\end{align}

\emph{Step 3: Global Idempotency of $M$.}
For any set $Z$ and any computation $mz \in M(Z)$, notice first that for any $m_{t\to mz} =
\mathrm{mApply}\;\eta_Z\;mz$, the monad associativity law~\eqref{eq:monad-assoc} and definition
\eqref{eq:mgz-to-mmz} yield
\begin{align}
  &\mathrm{mgz\!\to\!mmz}\bigl(\mathrm{mApply}\;\eta_Z\;mz\bigr) \nonumber \\
  &\quad= \mathrm{mApply}\;\bigl(t2mz \mapsto \mathrm{mPure}_{M(Z)}(t2mz(\ttstar))\bigr)\;\bigl(\mathrm{mApply}\;\eta_Z\;mz\bigr) \nonumber \\
  &\quad= \mathrm{mApply}\;\bigl(z \mapsto \mathrm{mApply}\;(t2mz \mapsto \mathrm{mPure}_{M(Z)}(t2mz(\ttstar)))\;(\eta_Z(z))\bigr)\;mz \nonumber \\
  &\quad= \mathrm{mApply}\;\bigl(z \mapsto \mathrm{mgz\!\to\!mmz}(\eta_Z(z))\bigr)\;mz. \label{eq:mgz-assoc}
\end{align}
Therefore, we obtain the complete equational chain:
\begin{align*}
  \mathrm{mMap}\;\mathrm{mPure}_Z\;mz
  &= \mathrm{mApply}\;\bigl(z \mapsto \mathrm{mPure}_{M(Z)}(\mathrm{mPure}_Z(z))\bigr)\;mz
    &&\text{(by \eqref{eq:mmap-mjoin})} \\
  &= \mathrm{mApply}\;\bigl(z \mapsto \mathrm{mgz\!\to\!mmz}(\eta_Z(z))\bigr)\;mz
    &&\text{(by \eqref{eq:idempotent-lemma} and ext.)} \\
  &= \mathrm{mgz\!\to\!mmz}\bigl(\mathrm{mApply}\;\eta_Z\;mz\bigr)
    &&\text{(by \eqref{eq:mgz-assoc})} \\
  &= \mathrm{mgz\!\to\!mmz}\bigl(\mathrm{mPure}_{G(Z)}(\_ \mapsto mz)\bigr)
    &&\text{(by \eqref{eq:eval-lemma})} \\
  &= \mathrm{mPure}_{M(Z)}(mz)
    &&\text{(by \eqref{eq:mgz-on-pure}, $h = (\_ \mapsto mz)$)}.
\end{align*}
Thus $\mathrm{mMap}\;\mathrm{mPure}_Z\;mz = \mathrm{mPure}_{M(Z)}(mz)$ for all $Z$ and $mz \in M(Z)$,
proving $\mathrm{Idempotent}(M)$.

\textbf{Direction ($\Longleftarrow$): $\mathrm{Idempotent}(M) \implies \mathrm{BasicFunctionalCategory}(\K(M), \KFF)$.}
Conversely, assume $\mathrm{Idempotent}(M)$ holds. For every set $Z \in \mathrm{ob}(\Set_\ell)$ and
$mz \in M(Z)$, define $mz\!\to\!mgmz \coloneqq \bigl(m \mapsto \mathrm{mPure}_{G(Z)}(\_ \mapsto m)\bigr)
: M(Z) \to M(G(Z))$. By the left unit law~\eqref{eq:monad-left-unit}, for each $z \in Z$ we have
$\mathrm{mPure}_{G(Z)}(\_ \mapsto \mathrm{mPure}_Z(z)) =
\mathrm{mApply}\;mz\!\to\!mgmz\;\bigl(\mathrm{mPure}_{M(Z)}(\mathrm{mPure}_Z(z))\bigr)$. Applying the
associativity law~\eqref{eq:monad-assoc}, definition~\eqref{eq:mmap-mjoin} of $\mathrm{mMap}$, the
idempotency hypothesis~\eqref{eq:idempotent}, and the left unit law~\eqref{eq:monad-left-unit} gives:
\begin{align}
  &\mathrm{mApply}\;\bigl(z \mapsto \mathrm{mPure}_{G(Z)}(\_ \mapsto \mathrm{mPure}_Z(z))\bigr)\;mz \nonumber \\
  &\quad= \mathrm{mApply}\;\bigl(z \mapsto \mathrm{mApply}\;mz\!\to\!mgmz\;\bigl(\mathrm{mPure}_{M(Z)}(\mathrm{mPure}_Z(z))\bigr)\bigr)\;mz \nonumber \\
  &\quad= \mathrm{mApply}\;mz\!\to\!mgmz\;\bigl(\mathrm{mApply}\;(z \mapsto \mathrm{mPure}_{M(Z)}(\mathrm{mPure}_Z(z)))\;mz\bigr) \nonumber \\
  &\quad= \mathrm{mApply}\;mz\!\to\!mgmz\;\bigl(\mathrm{mMap}\;\mathrm{mPure}_Z\;mz\bigr) \nonumber \\
  &\quad= \mathrm{mApply}\;mz\!\to\!mgmz\;\bigl(\mathrm{mPure}_{M(Z)}(mz)\bigr)
  \;=\; \mathrm{mPure}_{G(Z)}(\_ \mapsto mz). \label{eq:mapply-lemma}
\end{align}
Now define the candidate unit $\eta_Y : Y \to M(G(Y))$ for each $Y \in \mathrm{ob}(\K(M))$ by
\[
  \eta_Y(y) \;\coloneqq\; \mathrm{mPure}_{G(Y)}\bigl(\_ \mapsto \mathrm{mPure}_Y(y)\bigr).
\]
For any Kleisli morphism $y\!\to\!mx : Y \to M(X)$ and any $y \in Y$, applying~\eqref{eq:kleisli-comp},
\eqref{eq:mapply-lemma} at $mz = y\!\to\!mx\;(y)$, the left unit law~\eqref{eq:monad-left-unit} twice,
and~\eqref{eq:kleisli-G-map} yields naturality:
\begin{align*}
  &\bigl((y\!\to\!mx) \starK \eta_X\bigr)(y) \\
  &\quad= \mathrm{mApply}\;\bigl(x \mapsto \mathrm{mPure}_{G(X)}(\_ \mapsto \mathrm{mPure}_X(x))\bigr)\;\bigl(y\!\to\!mx\;(y)\bigr) \\
  &\quad= \mathrm{mPure}_{G(X)}\bigl(\_ \mapsto y\!\to\!mx\;(y)\bigr)
   \;=\; \mathrm{mPure}_{G(X)}\bigl(\_ \mapsto \mathrm{mApply}\;y\!\to\!mx\;(\mathrm{mPure}_Y(y))\bigr) \\
  &\quad= \mathrm{mApply}\;\bigl(gmy \mapsto \mathrm{mPure}_{G(X)}(t \mapsto \mathrm{mApply}\;y\!\to\!mx\;(gmy(t)))\bigr)\;\bigl(\mathrm{mPure}_{G(Y)}(\_ \mapsto \mathrm{mPure}_Y(y))\bigr) \\
  &\quad= \bigl(\eta_Y \starK G(y\!\to\!mx)\bigr)(y).
\end{align*}
Finally, by~\eqref{eq:kleisli-kff-g}, $\eta_Z(z) = \mathrm{mPure}_{G(Z)}(\_ \mapsto \mathrm{mPure}_Z(z))
= \KFF(g_Z)(z)$ holds definitionally by reflexivity for every set $Z$, establishing~\eqref{eq:eta-eq}.
\end{proof}

\subsection{The Terminal Subsingleton Theorem}

An immediate equational consequence of monad idempotency (which we derive directly here from
$\mathrm{IsPure}(M)$) is that any monadic computation over the terminal type $T$ must
be equal to $\mathrm{mPure}_T(\ttstar)$:

\begin{theorem}\label{thm:terminal-subsingleton}
Let $M$ be a monad on $\Set_\ell$ satisfying $\mathrm{IsPure}(M)$. Then $M$ satisfies the terminal
subsingleton property $\mathrm{MonadicTerminalTypeSubSingleton}(M)$: for every terminal monadic
computation $mt \in M(T)$,
\begin{equation}\label{eq:terminal-subsingleton}
  mt \;=\; \mathrm{mPure}_T(\ttstar).
\end{equation}
In particular, $M(T)$ is a subsingleton ($\forall m_1, m_2 \in M(T),\; m_1 = m_2$).
\end{theorem}

\begin{proof}
Define $\mathrm{mmt\!\to\!mt} \coloneqq \mathrm{mApply}\;(\_ \mapsto \mathrm{mPure}_T(\ttstar)) : M(M(T)) \to M(T)$.
Since $T = \{\ttstar\}$ has a unique inhabitant, $\mathrm{mPure}_T = (\_ \mapsto \mathrm{mPure}_T(\ttstar))$
by function extensionality, and for every $t \in T$ the left unit law~\eqref{eq:monad-left-unit} gives
$\mathrm{mmt\!\to\!mt}\bigl((\mathrm{mPure}_{M(T)} \circ \mathrm{mPure}_T)(t)\bigr) = \mathrm{mPure}_T(\ttstar)$.
Applying the right unit law~\eqref{eq:monad-right-unit}, associativity~\eqref{eq:monad-assoc}, the
purity law~\eqref{eq:ispure}, and the left unit law~\eqref{eq:monad-left-unit} yields:
\begin{align*}
  mt
  &= \mathrm{mApply}\;\mathrm{mPure}_T\;mt
    &&\text{(by \eqref{eq:monad-right-unit})} \\
  &= \mathrm{mApply}\;\bigl(\_ \mapsto \mathrm{mPure}_T(\ttstar)\bigr)\;mt
    &&\text{($T = \{\ttstar\}$)} \\
  &= \mathrm{mApply}\;\bigl(t \mapsto \mathrm{mmt\!\to\!mt}((\mathrm{mPure}_{M(T)} \circ \mathrm{mPure}_T)(t))\bigr)\;mt
    &&\text{(by \eqref{eq:monad-left-unit})} \\
  &= \mathrm{mmt\!\to\!mt}\bigl(\mathrm{mApply}\;(\mathrm{mPure}_{M(T)} \circ \mathrm{mPure}_T)\;mt\bigr)
    &&\text{(by \eqref{eq:monad-assoc})} \\
  &= \mathrm{mmt\!\to\!mt}\bigl(\mathrm{mPure}_{M(T)}(mt)\bigr)
    &&\text{(by \eqref{eq:ispure})} \\
  &= \mathrm{mPure}_T(\ttstar)
    &&\text{(by \eqref{eq:monad-left-unit})}.
  \qedhere
\end{align*}
\end{proof}

In my opinion, the proof above is an example of a \emph{theorem proving pearl} of
exceptional beauty. Using other propositions it proves that any $mt$ necessarily equals
$\mathrm{mPure}_T(\ttstar)$. 

\section{Classification of Concrete Computational Effects}\label{sec:examples}

We now evaluate $\mathrm{IsPure}(M)$ across seven standard computational effects, giving full
equational derivations for each case and summarizing the classification in Table~\ref{tab:effects}.

\begin{theorem}\label{thm:examples}
Let $\Set_\ell$ be a universe of sets, and for a set $A \in \mathrm{ob}(\Set_\ell)$ write
$\mathrm{Subsingleton}(A) \coloneqq (\forall a_1, a_2 \in A,\; a_1 = a_2)$.
\begin{enumerate}[label=(\arabic*),nosep]
  \item \textbf{Partiality ($\mathrm{Option}$, with $\mathrm{mPure}(z) = \mathrm{some}(z)$ and
  $\mathrm{mApply}\;f\;\mathrm{none} = \mathrm{none}$, $\mathrm{mApply}\;f\;(\mathrm{some}(z)) = f(z)$):}
  \[
    \neg\,\mathrm{IsPure}(\mathrm{Option}).
  \]
  \emph{Proof:} Choose the witness computation $mz = \mathrm{none} \in \mathrm{Option}(T)$. Then the
  left-hand side and right-hand side of the purity equation~\eqref{eq:ispure} evaluate to:
  \[
    \mathrm{mApply}\;(\mathrm{mPure} \circ \mathrm{mPure})\;\mathrm{none} \;=\; \mathrm{none},
    \qquad
    \mathrm{mPure}(\mathrm{none}) \;=\; \mathrm{some}(\mathrm{none}).
  \]
  Thus~\eqref{eq:ispure} forces $\mathrm{none} = \mathrm{some}(\mathrm{none})$ in
  $\mathrm{Option}(\mathrm{Option}(T))$, which is impossible since $\mathrm{none}$ and $\mathrm{some}$
  are disjoint constructors.
  \vspace{4pt}
  \item \textbf{Nondeterminism ($\mathrm{List}$, with $\mathrm{mPure}(z) = [z]$ and
  $\mathrm{mApply}\;f\;zs = \mathrm{concatMap}\;f\;zs$):}
  \[
    \neg\,\mathrm{IsPure}(\mathrm{List}).
  \]
  \emph{Proof:} Choose the two-element list $mz = [\ttstar, \ttstar] \in \mathrm{List}(T)$. Since
  $(\mathrm{mPure} \circ \mathrm{mPure})(z) = [[z]]$, the left-hand side and right-hand side of
  \eqref{eq:ispure} evaluate to:
  \[
    \mathrm{mApply}\;(\mathrm{mPure} \circ \mathrm{mPure})\;[\ttstar, \ttstar]
    \;=\; \bigl[[\ttstar],\, [\ttstar]\bigr],
    \qquad
    \mathrm{mPure}\bigl([\ttstar, \ttstar]\bigr)
    \;=\; \bigl[[\ttstar, \ttstar]\bigr].
  \]
  Equating the two lists in $\mathrm{List}(\mathrm{List}(T))$ and taking the tail (or comparing lengths
  $2 \neq 1$) yields $[[\ttstar]] = []$, a contradiction.
  \vspace{4pt}
  \item \textbf{Environment reading ($\mathrm{Read}\;R$, where $\mathrm{Read}\;R(Z) = (R \to Z)$,
  $\mathrm{mPure}(z) = (\_ \mapsto z)$, and $(\mathrm{mApply}\;f\;m)(r) = f(m(r))(r)$):}
  \[
    \mathrm{IsPure}(\mathrm{Read}\;R) \;\implies\; \mathrm{Subsingleton}(R).
  \]
  \emph{Proof:} Choose the identity reader $mz = \id_R = (r \mapsto r) \in \mathrm{Read}\;R(R)$. For
  any $z \in R$, $(\mathrm{mPure} \circ \mathrm{mPure})(z) = (\_ \mapsto (\_ \mapsto z)) : R \to (R \to R)$.
  Evaluating both sides of~\eqref{eq:ispure} on $mz = \id_R$ gives:
  \begin{align*}
    \mathrm{LHS}
    &= \mathrm{mApply}\;(\mathrm{mPure} \circ \mathrm{mPure})\;\id_R \\
    &= \bigl(r_1 \mapsto ((\mathrm{mPure} \circ \mathrm{mPure})(\id_R(r_1)))(r_1)\bigr)
     \;=\; \bigl(r_1 \mapsto (r_2 \mapsto r_1)\bigr), \\
    \mathrm{RHS}
    &= \mathrm{mPure}(\id_R)
     \;=\; \bigl(r_1 \mapsto \id_R\bigr)
     \;=\; \bigl(r_1 \mapsto (r_2 \mapsto r_2)\bigr).
  \end{align*}
  For any $r_1, r_2 \in R$, applying $\mathrm{LHS} = \mathrm{RHS}$ to $r_1$ and then to $r_2$ gives
  $r_1 = \mathrm{LHS}(r_1)(r_2) = \mathrm{RHS}(r_1)(r_2) = r_2$.
  \vspace{4pt}
  \item \textbf{Logging ($\mathrm{Write}\;W$, where $\mathrm{Write}\;W(Z) = Z \times W$ for a monoid
  $(W, e, \bullet)$, $\mathrm{mPure}(z) = (z, e)$, and $\mathrm{mApply}\;f\;(z, w_0) = (y, w_0 \bullet w_1)$
  where $(y, w_1) = f(z)$):}
  \[
    \mathrm{IsPure}(\mathrm{Write}\;W) \;\implies\; \mathrm{Subsingleton}(W).
  \]
  \emph{Proof:} Let $w \in W$ and consider the witness $mz = (e, w) \in \mathrm{Write}\;W(W)$. Since
  $(\mathrm{mPure} \circ \mathrm{mPure})(z) = ((z, e), e)$, the two sides of~\eqref{eq:ispure} at
  $mz = (e, w)$ are:
  \begin{align*}
    \mathrm{LHS} &= \mathrm{mApply}\;(\mathrm{mPure} \circ \mathrm{mPure})\;(e, w) \;=\; \bigl((e, e),\, w \bullet e\bigr), \\
    \mathrm{RHS} &= \mathrm{mPure}\bigl((e, w)\bigr) \;=\; \bigl((e, w),\, e\bigr).
  \end{align*}
  Taking the first projection $\pi_1$ of $\mathrm{LHS} = \mathrm{RHS}$ gives $(e, e) = (e, w)$, and
  taking the second projection $\pi_2$ gives $e = w$. Hence for any $w_1, w_2 \in W$, $w_1 = e = w_2$.
  \vspace{4pt}
  \item \textbf{Mutable state ($\mathrm{State}\;S$, where $\mathrm{State}\;S(Z) = (S \to Z \times S)$,
  $\mathrm{mPure}(z) = (s \mapsto (z, s))$, and
  $(\mathrm{mApply}\;f\;m)(s) = f(z)(s')$ where $(z, s') = m(s)$):}
  \[
    \mathrm{IsPure}(\mathrm{State}\;S) \;\implies\; \mathrm{Subsingleton}(S).
  \]
  \emph{Proof:} Choose the state-reading computation $mz = \bigl(s \mapsto (s, s)\bigr) \in \mathrm{State}\;S(S)$.
  For any $z \in S$, we have
  \[
    (\mathrm{mPure} \circ \mathrm{mPure})(z)
    \;=\; \mathrm{mPure}\bigl(s' \mapsto (z, s')\bigr)
    \;=\; \Bigl(s \mapsto \bigl((s' \mapsto (z, s')),\, s\bigr)\Bigr).
  \]
  Therefore, the left-hand side of~\eqref{eq:ispure} applied to $mz = (s \mapsto (s, s))$ is:
  \begin{align*}
    \mathrm{LHS}(s_1)
    &= \bigl(\mathrm{mApply}\;(\mathrm{mPure} \circ \mathrm{mPure})\;mz\bigr)(s_1) \\
    &= \bigl((\mathrm{mPure} \circ \mathrm{mPure})(s_1)\bigr)(s_1)
    \;=\; \bigl((s' \mapsto (s_1, s')),\, s_1\bigr),
  \end{align*}
  while the right-hand side of~\eqref{eq:ispure} applied to $mz = (s \mapsto (s, s))$ is:
  \[
    \mathrm{RHS}(s_1)
    \;=\; \bigl(\mathrm{mPure}(mz)\bigr)(s_1)
    \;=\; (mz,\, s_1)
    \;=\; \bigl((s' \mapsto (s', s')),\, s_1\bigr).
  \]
  For any $s_1, s_2 \in S$, evaluating $\mathrm{LHS} = \mathrm{RHS}$ at $s_1$ and taking the first
  projection $\pi_1$ yields the equality of functions $(s' \mapsto (s_1, s')) = (s' \mapsto (s', s'))$
  in $S \to S \times S$. Evaluating both functions at $s' = s_2$ gives $(s_1, s_2) = (s_2, s_2)$, and
  taking the first projection $\pi_1$ yields $s_1 = s_2$.
  \vspace{4pt}
  \item \textbf{Continuations ($\mathrm{Cont}\;R$, where $\mathrm{Cont}\;R(Z) = ((Z \to R) \to R)$,
  $\mathrm{mPure}(z) = (k \mapsto k(z))$, and
  $(\mathrm{mApply}\;f\;m)(k) = m(z \mapsto f(z)(k))$):}
  \[
    \mathrm{IsPure}(\mathrm{Cont}\;R) \;\implies\; \mathrm{Subsingleton}(R).
  \]
  \emph{Proof:} Let $r_1, r_2 \in R$, and consider the constant continuation computation
  $m_1 \coloneqq (\_ \mapsto r_1) \in \mathrm{Cont}\;R(R) = ((R \to R) \to R)$, which ignores its
  continuation argument $k : R \to R$ and returns $r_1$. In $\mathrm{Cont}\;R(\mathrm{Cont}\;R(R))$,
  an outer continuation is a function $K : \mathrm{Cont}\;R(R) \to R$. The left-hand side of
  \eqref{eq:ispure} at $mz = m_1$, applied to any outer continuation $K : \mathrm{Cont}\;R(R) \to R$,
  evaluates to:
  \begin{align*}
    \mathrm{LHS}(K)
    &= \bigl(\mathrm{mApply}\;(\mathrm{mPure} \circ \mathrm{mPure})\;m_1\bigr)(K) \\
    &= m_1\bigl(z \mapsto ((\mathrm{mPure} \circ \mathrm{mPure})(z))(K)\bigr) \;=\; r_1,
  \end{align*}
  because $m_1 = (\_ \mapsto r_1)$ returns $r_1$ regardless of its argument. Conversely, the right-hand
  side of~\eqref{eq:ispure} at $mz = m_1$, applied to $K$, evaluates to:
  \[
    \mathrm{RHS}(K)
    \;=\; \bigl(\mathrm{mPure}(m_1)\bigr)(K)
    \;=\; K(m_1).
  \]
  Now choose the constant outer continuation $K_2 \coloneqq (\_ \mapsto r_2) : \mathrm{Cont}\;R(R) \to R$.
  Evaluating $\mathrm{LHS} = \mathrm{RHS}$ at $K = K_2$ gives
  \[
    r_1 \;=\; \mathrm{LHS}(K_2) \;=\; \mathrm{RHS}(K_2) \;=\; K_2(m_1) \;=\; r_2.
  \]
  \vspace{4pt}
  \item \textbf{Input/Output and Exceptions ($\mathrm{ErrorState}\;E\;S$, where
  $\mathrm{ErrorState}\;E\;S(Z) = (S \to (E \times S) \uplus (Z \times S))$,
  $\mathrm{mPure}(z) = (s \mapsto \mathrm{inr}(z, s))$, and $\mathrm{mApply}\;f\;m$ propagates
  $\mathrm{inl}(e, s')$ and applies $f(z)(s')$ on $\mathrm{inr}(z, s')$):}
  \[
    \forall e \in E,\; \forall s \in S,\qquad \neg\,\mathrm{IsPure}(\mathrm{ErrorState}\;E\;S).
  \]
  \emph{Proof:} Given $e \in E$ and $s \in S$, define the failing computation
  $m \coloneqq \bigl(s' \mapsto \mathrm{inl}(e, s')\bigr) \in \mathrm{ErrorState}\;E\;S(T)$. Evaluating
  both sides of~\eqref{eq:ispure} on $m$ and then at initial state $s \in S$ yields:
  \[
    \bigl(\mathrm{mApply}\;(\mathrm{mPure} \circ \mathrm{mPure})\;m\bigr)(s) \;=\; \mathrm{inl}(e, s),
    \qquad
    \bigl(\mathrm{mPure}(m)\bigr)(s) \;=\; \mathrm{inr}(m, s),
  \]
  contradicting the disjointness of $\mathrm{inl}$ and $\mathrm{inr}$.
\end{enumerate}
\end{theorem}

\begin{table}[htbp]
\centering
\scriptsize
\setlength{\tabcolsep}{3.5pt}
\renewcommand{\arraystretch}{1.25}
\begin{tabular}{@{}llllll@{}}
\toprule
\textbf{Monad $M$} & \textbf{Carrier $M(Z)$} & \textbf{Terminal $M(T)$} & \textbf{Witness $mz$} & \textbf{LHS vs.\ RHS of~\eqref{eq:ispure}} & \textbf{$\mathrm{IsPure}(M)$} \\
\midrule
$\mathrm{Option}$ & $1 + Z$ & $1 + 1$ & $\mathrm{none}$ & $\mathrm{none} \neq \mathrm{some}(\mathrm{none})$ & False \\
$\mathrm{List}$ & $\mathrm{List}(Z)$ & $\mathrm{List}(1) \cong \mathbb{N}$ & $[\ttstar, \ttstar]$ & $[[\ttstar],[\ttstar]] \neq [[\ttstar,\ttstar]]$ & False \\
$\mathrm{Read}\;R$ & $R \to Z$ & $R \to 1 \cong 1$ & $\id_R$ & $(r_1, r_2 \mapsto r_1) = (r_1, r_2 \mapsto r_2)$ & $\iff |R| \le 1$ \\
$\mathrm{Write}\;W$ & $Z \times W$ & $1 \times W \cong W$ & $(e, w)$ & $((e,e), w) = ((e,w), e)$ & $\iff |W| \le 1$ \\
$\mathrm{State}\;S$ & $S \to Z \times S$ & $S \to 1 \times S \cong S^S$ & $s \mapsto (s, s)$ & $((s' \mapsto (s_1,s')), s_1) = (mz, s_1)$ & $\iff |S| \le 1$ \\
$\mathrm{Cont}\;R$ & $(Z \to R) \to R$ & $(1 \to R) \to R \cong R^R$ & $\_ \mapsto r_1$ & $(\_ \mapsto r_1) = (K \mapsto K(m_1))$ & $\iff |R| \le 1$ \\
$\mathrm{ErrorState}\;E\;S$ & $S \to (E + Z) \times S$ & $S \to (E + 1) \times S$ & $s' \mapsto \mathrm{inl}(e, s')$ & $\mathrm{inl}(e, s) \neq \mathrm{inr}(m, s)$ & False ($E, S \neq \emptyset$) \\
\bottomrule
\end{tabular}
\caption{Summary of the seven standard computational effects, their terminal carriers $M(T)$, their
purity witness computations, and their exact $\mathrm{IsPure}(M)$ classification (where $|A| \le 1$
denotes $\mathrm{Subsingleton}(A)$).}
\label{tab:effects}
\end{table}

\subsection{Strict Separation Between Basic Functional Category Purity and Terminal Subsingleton Triviality}

As Table~\ref{tab:effects} illustrates, the terminal subsingleton condition
$\mathrm{MonadicTerminalTypeSubSingleton}(M)$ ($M(T) \cong T$, Theorem~\ref{thm:terminal-subsingleton})
suffices to rule out $\mathrm{Option}$, $\mathrm{List}$, $\mathrm{Write}\;W$ ($|W| > 1$),
$\mathrm{State}\;S$ ($|S| > 1$), $\mathrm{Cont}\;R$ ($|R| > 1$), and $\mathrm{ErrorState}\;E\;S$
($E, S \neq \emptyset$). However, $\mathrm{Read}\;R$ satisfies
\[
  \mathrm{Read}\;R(T) \;=\; (R \to T) \;\cong\; T
\]
for \emph{every} environment set $R$. We now formulate and prove a formal separation theorem showing
that the Basic Functional Category condition on $(\K(M), \KFF)$ (and in particular the existence of a
natural transformation $\eta : \Id_{\K(M)} \Rightarrow G$ satisfying the terminal unit
law~\eqref{eq:eta-eq-T}) is \textbf{strictly stronger} than the terminal subsingleton property
$\mathrm{MonadicTerminalTypeSubSingleton}(M)$.

\begin{theorem}[Strict separation theorem]\label{thm:strict-separation}
For monads $M$ on $\Set_\ell$:
\begin{enumerate}[label=(\roman*),nosep]
  \item \textbf{Universal Implication:} For every monad $M$ on $\Set_\ell$,
  \[
    \mathrm{BasicFunctionalCategory}\bigl(\K(M),\, \KFF\bigr)
    \;\implies\;
    \mathrm{MonadicTerminalTypeSubSingleton}(M).
  \]
\vspace{4 pt}
  \item \textbf{Strictness (Failure of the Converse):} For every set $R \in \mathrm{ob}(\Set_\ell)$
  containing at least two distinct elements $r_1 \neq r_2$ (for instance, $R = \{0, 1\}$), the reader
  monad $M = \mathrm{Read}\;R$ satisfies:
  \begin{enumerate}[label=(\alph*),nosep]
    \item $\mathrm{MonadicTerminalTypeSubSingleton}(\mathrm{Read}\;R)$ \textbf{holds} ($\forall mt \in
    \mathrm{Read}\;R(T),\; mt = \mathrm{mPure}_T(\ttstar)$);
\vspace{4 pt}
    \item $\mathrm{BasicFunctionalCategory}\bigl(\K(\mathrm{Read}\;R),\, \KFF\bigr)$ \textbf{fails}
    (no natural transformation $\eta : \Id_{\K(\mathrm{Read}\;R)} \Rightarrow G$ satisfies the terminal
    unit law~\eqref{eq:eta-eq-T}).
  \end{enumerate}
\end{enumerate}
\end{theorem}

\begin{proof}
\textbf{Part (i):} Suppose $(\K(M), \KFF)$ is a Basic Functional Category. By
Theorem~\ref{thm:kleisli-purity}, $M$ satisfies $\mathrm{Idempotent}(M)$, which by
Lemma~\ref{lem:idem-ispure} is equivalent to $\mathrm{IsPure}(M)$. Applying
Theorem~\ref{thm:terminal-subsingleton} yields $\mathrm{MonadicTerminalTypeSubSingleton}(M)$.

\textbf{Part (ii)(a):} Let $mt \in \mathrm{Read}\;R(T) = (R \to T)$. For every $r \in R$, since
$T = \{\ttstar\}$ has a unique inhabitant, $mt(r) = \ttstar = (\_ \mapsto \ttstar)(r) =
\mathrm{mPure}_T(\ttstar)(r)$. By function extensionality, $mt = \mathrm{mPure}_T(\ttstar)$, so
$\mathrm{MonadicTerminalTypeSubSingleton}(\mathrm{Read}\;R)$ holds for every set $R$.

\textbf{Part (ii)(b):} Suppose for contradiction that $(\K(\mathrm{Read}\;R), \KFF)$ is a Basic
Functional Category (or merely that there exists a natural transformation
$\eta : \Id_{\K(\mathrm{Read}\;R)} \Rightarrow G$ satisfying the terminal unit law~\eqref{eq:eta-eq-T}).
By Theorem~\ref{thm:kleisli-purity}(i) and Lemma~\ref{lem:idem-ispure}, $\mathrm{IsPure}(\mathrm{Read}\;R)$
must hold. By Theorem~\ref{thm:examples}(3), $\mathrm{IsPure}(\mathrm{Read}\;R)$ implies
$\mathrm{Subsingleton}(R)$, forcing $r_1 = r_2$, which contradicts $r_1 \neq r_2$.

To see this failure directly on the explicit separating model $M = \mathrm{Read}\;R$ with two distinct
environment values $r_1 \neq r_2 \in R$ (such as the two-element boolean environment $R = \{0, 1\}$ with
$r_1 = 0$ and $r_2 = 1$), note that in $\K(\mathrm{Read}\;R)$, the candidate unit forced by~\eqref{eq:eta-eq}
is $\eta_Z(z) = \bigl(r_{\mathrm{outer}} \mapsto (\_ \mapsto (r_{\mathrm{inner}} \mapsto z))\bigr)$, and
at $Z = T$ the terminal law~\eqref{eq:eta-eq-T} forces
$\eta_T(\ttstar) = \bigl(r_{\mathrm{outer}} \mapsto (\_ \mapsto (r_{\mathrm{inner}} \mapsto \ttstar))\bigr)$.
Consider the environment-reading Kleisli morphism
\[
  f \;\coloneqq\; (\_ \mapsto \id_R) \;=\; \bigl(t \mapsto (r \mapsto r)\bigr)
  \;\in\; \Hom_{\K(\mathrm{Read}\;R)}(T, R).
\]
Evaluating the naturality equation $f \starK \eta_R = \eta_T \starK G(f)$ at $\ttstar \in T$, outer
environment $r_1 \in R$, dummy terminal argument $\ttstar \in T$, and inner environment $r_2 \in R$ gives:
\begin{align*}
  \bigl((f \starK \eta_R)(\ttstar)(r_1)(\ttstar)\bigr)(r_2)
  &= \bigl(\eta_R(f(\ttstar)(r_1))(r_1)(\ttstar)\bigr)(r_2)
   \;=\; \bigl(\eta_R(r_1)(r_1)(\ttstar)\bigr)(r_2) \;=\; r_1, \\
  \bigl((\eta_T \starK G(f))(\ttstar)(r_1)(\ttstar)\bigr)(r_2)
  &= \bigl(G(f)(\eta_T(\ttstar)(r_1))(r_1)(\ttstar)\bigr)(r_2) \\
  &= f\bigl(\eta_T(\ttstar)(r_1)(\ttstar)(r_2)\bigr)(r_2)
   \;=\; f(\ttstar)(r_2) \;=\; r_2.
\end{align*}
On the left-hand side, binding $f$ before $\eta_R$ reads the \emph{outer} environment $r_1$ and freezes
it inside the thunk; on the right-hand side, postcomposing $f$ inside $G(f)$ defers reading the
environment until the \emph{inner} evaluation at $r_2$. Since $r_1 \neq r_2$, naturality of $\eta$ fails
for $\mathrm{Read}\;R$ whenever $|R| \ge 2$, even though $\mathrm{Read}\;R(T) = (R \to T) \cong T$ is a
subsingleton.
\end{proof}

\subsection{Scope, Implications for Commutative Effects, and Limitations of the Framework}\label{subsec:limitations}

Our results establish a sharp boundary for the pointfree Basic Functional Category and CoEndoYoneda
framework, which clarifies both its power as a purity specification and its scope for modeling effects:

\begin{enumerate}[label=(\arabic*),nosep]
  \item \textbf{Exclusion of non-idempotent and commutative effects:}
  By Theorem~\ref{thm:kleisli-purity}, for a Kleisli category $\K(M)$ equipped with the canonical
  identity-on-objects embedding $\KFF : \Set_\ell \to \K(M)$, the Basic Functional Category structure
  (and hence the CoEndoYoneda Isomorphism Theorem~\ref{thm:coendoyoneda}) holds if and only if $M$ is an
  idempotent monad. Consequently, the CoEndoYoneda isomorphism with respect to $\KFF$ does \emph{not}
  apply to Kleisli categories of non-idempotent monads such as $\mathrm{Option}$, $\mathrm{List}$,
  $\mathrm{State}\;S$ ($|S| \ge 2$), $\mathrm{Cont}\;R$ ($|R| \ge 2$), or $\mathrm{ErrorState}\;E\;S$
  ($E, S \neq \emptyset$). Moreover, unlike premonoidal centrality~\cite{PowerRobinson1997,LevyPowerThielecke2003}---which
  requires effect commutativity $(f \ltimes h) = (f \rtimes h)$ and therefore includes commutative
  effects such as $\mathrm{Read}\;R$ and commutative $\mathrm{Write}\;W$ in the center---our condition
  requires that \emph{every} computation commute with the thunking endofunctor $G(Z) = (T \to M(Z))$ via
  $\eta$. As shown in Theorem~\ref{thm:strict-separation}, for $\mathrm{Read}\;R$ ($|R| \ge 2$),
  binding an environment query before $\eta$ freezes the outer environment $r_1$, whereas pushing the
  query into $G$ reads the inner environment $r_2$; for $\mathrm{Write}\;W$ ($|W| \ge 2$), binding a log
  emission before $\eta$ emits the log at the outer level, whereas pushing it into $G$ traps the log
  inside the thunk $T \to (Z \times W)$. Thus, the Basic Functional Category axiom characterizes
  \emph{strict effectlessness} (referential transparency across thunking) rather than effect
  commutativity.

  \vspace{3pt}
  \item \textbf{Which non-trivial monads are idempotent?}
  Although the seven effect families in Theorem~\ref{thm:examples} collapse to trivial parameter sets
  when idempotent, idempotent monads in general category theory are far from trivial: on a general
  category $\mathcal{C}$, idempotent monads classify \emph{reflective subcategories} (and localization /
  modality operators, such as open/closed modalities, sheafification, or propositional truncation in
  homotopy type theory)~\cite{Borceux1994,Kelly1980,nLabIdempotentMonad}. On $\Set_\ell$, where every
  non-empty set splits over $T$, the only reflective subcategories are the full category $\Set_\ell$
  ($M \cong \Id_{\Set_\ell}$) and the terminal subcategory ($M(Z) \cong T$, or $M(\emptyset) = \emptyset$
  and $M(Z) \cong T$ for $Z \neq \emptyset$). Thus, on $\Set_\ell$, our characterization accurately
  certifies that a monadic program category $(\K(M), \KFF)$ behaves purely if and only if $M$ introduces
  no computational side effects beyond pure values (or trivial termination).

  \vspace{3pt}
  \item \textbf{Alternative pointfree characterizations for benign or commutative effects:}
  If one wishes to relax strict purity to admit specific classes of benign effects in a pointfree
  setting, our hierarchy of theorems suggests three natural directions:
  \begin{enumerate}[label=(\alph*),nosep]
    \item \emph{Discardable / read-only effects} (such as $\mathrm{Read}\;R$): characterized by the
    strictly weaker terminal subsingleton law $M(T) \cong T$ (Theorem~\ref{thm:terminal-subsingleton}),
    which demands that computations returning no data carry no observable side-effect state, and corresponds
    categorically to the terminal object $\KFF(T)$ remaining terminal in $\K(M)$ (semicartesian / affine
    structure);
    \item \emph{Commutative effects} (such as $\mathrm{Read}\;R$ and commutative $\mathrm{Write}\;W$):
    characterized pointfreely by central morphisms in premonoidal and Freyd
    categories~\cite{PowerRobinson1997,LevyPowerThielecke2003};
    \item \emph{Restricted or relative CoEndoYoneda structures:} restricting the naturality of
    $\eta : \Id_{\mathcal{C}} \Rightarrow G$ to the subcategory of pure (or thunkable~\cite{Fuhrmann1999})
    morphisms inside an ambient effectful category $\mathcal{C}$, or varying the functional functor
    $\FF : \Set_\ell \to \mathcal{C}$ beyond the identity-on-objects Kleisli inclusion $\KFF$.
  \end{enumerate}
\end{enumerate}

\section{Formal Verification in Cubical Agda and Lean 4}\label{sec:formalization}

Both formalizations are 100\% machine-checked with zero postulates, zero unproved axioms, and zero
\texttt{sorry}s. Table~\ref{tab:alignment} provides the exact alignment between the mathematical results
in this paper, the Cubical Agda development (\texttt{CoEndoYoneda.lagda.md}~\cite{CubicalAgdaRepo}),
and the Palomar-verified Lean~4 formalization (\texttt{coEndoYoneda-and-purity}~\cite{PalomarLeanRepo}).

\begin{table}[htbp]
\centering
\scriptsize
\renewcommand{\arraystretch}{1.2}
\begin{tabular}{@{}llll@{}}
\toprule
\textbf{Paper Reference} & \textbf{Cubical Agda (\texttt{CoEndoYoneda.lagda.md})} & \textbf{Lean 4 (\texttt{CoEndoYonedaAndPurity})} & \textbf{Status} \\
\midrule
Definition~\ref{def:bfc} & \texttt{BasicFunctionalCategory} & \texttt{BasicFunctionalCategory} & Verified \\
Lemma~\ref{lem:g-id} & \texttt{lemma}, \texttt{lemma-T} & \texttt{lemma\_g}, \texttt{lemma\_T} & Proved \\
Proposition~\ref{prop:eta-eq-T} & \texttt{η-eq-T} & \texttt{η\_eq\_T} & Proved \\
Definition~\ref{def:fc} & \texttt{FunctionalCategory} & \texttt{FunctionalCategory} & Verified \\
Theorem~\ref{thm:coendoyoneda}(i) & \texttt{fz→τz\_commute} & \texttt{fz\_to\_τz\_commute} & Proved \\
Theorem~\ref{thm:coendoyoneda}(ii) & \texttt{to\_from\_proof} & \texttt{to\_from\_proof\_hom}, \texttt{to\_from\_proof} & Proved \\
Theorem~\ref{thm:coendoyoneda}(iii) & \texttt{from\_to\_proof'}, \texttt{from\_to\_proof} & \texttt{from\_to\_proof'}, \texttt{from\_to\_proof} & Proved \\
Theorem~\ref{thm:coendoyoneda} (Iso) & \texttt{coEndoYonedaIso} & \texttt{coEndoYonedaIso} & Proved \\
Proposition~\ref{prop:origin-eta} & \texttt{ηFFT-is-η-at-FFT} & \texttt{ηFFT\_is\_η\_at\_FFT} & Proved \\
Proposition~\ref{prop:set-fc} & \texttt{SetFunctionalCategory} & \texttt{set\_is\_functional\_category} & Proved \\
Lemma~\ref{lem:idem-ispure} & \texttt{idempotent-equiv-is-pure} & \texttt{idempotent\_iff\_isPure} & Proved \\
Theorem~\ref{thm:kleisli-purity} & \texttt{idempotent-equiv-kleisli-...} & \texttt{kleisliBasicFunctionalCategory\_iff\_...} & Proved \\
Theorem~\ref{thm:terminal-subsingleton} & \texttt{is-pure-implies-monadic-...} & \texttt{isPure\_implies\_monadicTerminalType...} & Proved \\
Theorem~\ref{thm:examples}(1)--(2) & \texttt{maybe-not-pure}, \texttt{list-not-pure} & \texttt{maybe\_not\_pure}, \texttt{list\_not\_pure} & Proved \\
Theorem~\ref{thm:examples}(3)--(6) & \texttt{read/write/state/cont-purity-implies} & \texttt{read/write/state/cont\_purity\_implies} & Proved \\
Theorem~\ref{thm:examples}(7) & \texttt{io-not-pure} & \texttt{io\_not\_pure} & Proved \\
\bottomrule
\end{tabular}
\caption{Alignment between the paper, the Cubical Agda formalization, and the Lean~4 Palomar formalization.}
\label{tab:alignment}
\end{table}

\begin{thebibliography}{99}
\raggedright

\bibitem{Borceux1994}
F. Borceux,
\emph{Handbook of Categorical Algebra 2: Categories and Structures},
Encyclopedia of Mathematics and its Applications, vol.~51, Cambridge University Press, 1994.

\bibitem{DuponcheelAcademia}
L. Duponcheel,
\emph{A pointfree Yoneda lemma for endofunctors of functional categories},
Working paper, \url{https://www.academia.edu/109957588}.

\bibitem{CubicalAgdaRepo}
L. Duponcheel,
\emph{Cubical Agda formalization: CoEndoYoneda and its relation to Purity}, \\
\url{https://github.com/LucDuponcheelAtGitHub/AgdaCoEndoYonedaAndPurity}, 2026.

\bibitem{PalomarLeanRepo}
L. Duponcheel,
\emph{Palomar Lean~4 formalization: CoEndoYoneda and Purity}, \\
\url{https://github.com/LucDuponcheelAtGitHub/coEndoYoneda-and-purity}, 2026.

\bibitem{PSBP}
L. Duponcheel,
\emph{Program Specification Based Programming}, \\
\url{https://github.com/LucDuponcheelAtGitHub/AgdaPSBP}, 2026.

\bibitem{Fuhrmann1999}
C. F\"{u}hrmann,
\emph{Direct models of the computational lambda-calculus},
Electronic Notes in Theoretical Computer Science \textbf{20} (1999), 245--292,
\url{https://doi.org/10.1016/S1571-0661(04)80078-1}.

\bibitem{Kelly1980}
G. M. Kelly,
\emph{A unified treatment of transfinite constructions for free algebras, free monoids, colimits,
associated sheaves, and so on},
Bulletin of the Australian Mathematical Society \textbf{22} (1980), no.~1, 1--83,
\url{https://doi.org/10.1017/S0004972700006353}.

\bibitem{Lawvere1963}
F. W. Lawvere,
\emph{Functorial semantics of algebraic theories},
Proceedings of the National Academy of Sciences \textbf{50} (1963), no.~5, 869--872,
\url{https://doi.org/10.1073/pnas.50.5.869}.

\bibitem{LevyPowerThielecke2003}
P. B. Levy, J. Power, and H. Thielecke,
\emph{Modelling environments in call-by-value programming languages},
Information and Computation \textbf{185} (2003), no.~2, 182--210,
\url{https://doi.org/10.1016/S0890-5401(03)00088-9}.

\bibitem{Moggi1991}
E. Moggi,
\emph{Notions of computation and monads},
Information and Computation \textbf{93} (1991), no.~1, 55--92,
\url{https://doi.org/10.1016/0890-5401(91)90052-4}.

\bibitem{nLabIdempotentMonad}
nLab authors,
\emph{Idempotent monad}, \\
\url{https://ncatlab.org/nlab/show/idempotent+monad}.

\bibitem{nLabWellPointed}
nLab authors,
\emph{Well-pointed endofunctor}, \\
\url{https://ncatlab.org/nlab/show/well-pointed+endofunctor}.

\bibitem{PlotkinPower2002}
G. Plotkin and J. Power,
\emph{Notions of computation determine monads},
in: \emph{Foundations of Software Science and Computation Structures (FoSSaCS 2002)},
Lecture Notes in Computer Science, vol.~2303, Springer, 2002, pp.~342--356,
\url{https://doi.org/10.1007/3-540-45931-6_24}.

\bibitem{PlotkinPower2003}
G. Plotkin and J. Power,
\emph{Algebraic operations and generic effects},
Applied Categorical Structures \textbf{11} (2003), no.~1, 69--94,
\url{https://doi.org/10.1023/A:1023064908962}.

\bibitem{PowerRobinson1997}
J. Power and E. Robinson,
\emph{Premonoidal categories and notions of computation},
Mathematical Structures in Computer Science \textbf{7} (1997), no.~5, 453--468,
\url{https://doi.org/10.1017/S0960129597002375}.

\bibitem{PowerThielecke1999}
J. Power and H. Thielecke,
\emph{Closed Freyd- and $\kappa$-categories},
in: \emph{Automata, Languages and Programming (ICALP 1999)},
Lecture Notes in Computer Science, vol.~1644, Springer, 1999, pp.~625--634,
\url{https://doi.org/10.1007/3-540-48523-6_59}.

\end{thebibliography}

\end{document}
